\documentclass[UTF8,12pt,a4paper]{ctexart}

% ==============================================================================
% 宏包加载与基础配置
% ==============================================================================
\usepackage[top=2.5cm, bottom=2.5cm, left=2.5cm, right=2.5cm]{geometry}
\usepackage{amsmath, amssymb, amsthm, mathtools}
\usepackage{booktabs, tabularx, longtable, multirow}
\usepackage{graphicx}
\usepackage{xcolor}
\usepackage{tikz}
\usetikzlibrary{arrows.meta, positioning, shapes.geometric, shapes.misc, calc, trees, backgrounds, fit}
\usepackage{pgfplots}
\pgfplotsset{compat=1.18}
\usepackage{tcolorbox}
\tcbuselibrary{skins, breakable}
\usepackage{enumitem}
\usepackage{fancyhdr}
\usepackage{lastpage}
\usepackage{titlesec}
\usepackage{caption}
\usepackage{subcaption}
\usepackage{cite}
\usepackage[hidelinks, colorlinks=true, linkcolor=navyblue, citecolor=navyblue, urlcolor=navyblue]{hyperref}
\usepackage{cleveref}

% ==============================================================================
% 色彩与版面定义
% ==============================================================================
\definecolor{navyblue}{RGB}{27, 54, 93}
\definecolor{slateblue}{RGB}{44, 82, 130}
\definecolor{charcoal}{RGB}{45, 55, 72}
\definecolor{darkgray}{RGB}{74, 85, 104}
\definecolor{lightbg}{RGB}{247, 250, 252}
\definecolor{bordercolor}{RGB}{226, 232, 240}
\definecolor{accentgreen}{RGB}{40, 116, 101}
\definecolor{accentamber}{RGB}{180, 83, 9}

\linespread{1.3} % 1.3 倍行距
\setlength{\parskip}{0.5em}
\setlength{\parindent}{2em}

% 标题格式定制
\titleformat{\section}{\color{navyblue}\Large\bfseries}{\thesection}{1em}{}
\titleformat{\subsection}{\color{slateblue}\large\bfseries}{\thesubsection}{1em}{}
\titleformat{\subsubsection}{\color{charcoal}\normalsize\bfseries}{\thesubsubsection}{1em}{}

% 页眉页脚配置
\setlength{\headheight}{15pt}
\pagestyle{fancy}
\fancyhf{}
\fancyhead[L]{\small\color{darkgray}企业大模型使用数据差分隐私分析系统技术方案}
\fancyhead[R]{\small\color{darkgray}Private LLM Analytics}
\fancyfoot[L]{\small\color{darkgray}技术方案版本: v1.1-dev}
\fancyfoot[R]{\small\color{darkgray}第 \thepage\ 页\,/\,共 \pageref{LastPage}\ 页}
\renewcommand{\headrulewidth}{0.5pt}
\renewcommand{\footrulewidth}{0.5pt}

% 自定义高亮框
\newtcolorbox{techbox}[2][]{
  enhanced,
  breakable,
  colback=lightbg,
  colframe=navyblue,
  coltitle=white,
  fonttitle=\bfseries\small,
  title={#2},
  boxrule=0.7pt,
  arc=2pt,
  left=8pt,
  right=8pt,
  top=6pt,
  bottom=6pt,
  #1
}

\newtcolorbox{definitionbox}[2][]{
  enhanced,
  breakable,
  colback=white,
  colframe=slateblue,
  coltitle=white,
  fonttitle=\bfseries\small,
  title={定义: #2},
  boxrule=0.7pt,
  arc=2pt,
  left=8pt,
  right=8pt,
  top=6pt,
  bottom=6pt,
  #1
}

% ==============================================================================
% 封面与元数据
% ==============================================================================
\title{
  \vspace{-1.5cm}
  {\Huge\bfseries\color{navyblue} 企业大模型使用数据\\[0.35cm]差分隐私分析系统技术方案}\\[0.6cm]
  {\large\bfseries\color{slateblue} Private LLM Analytics: Technical Architecture and Privacy-Preserving Analytics Design}\\[0.4cm]
  {\normalsize\color{darkgray} 系统架构设计、差分隐私机制证明、全生命周期记账与工程实现规范}
}

\author{
  \textbf{系统架构与数据安全联合研发团队} \\
  \textit{Differential Privacy Engineering Group}
}

\date{
  \vspace{0.4cm}
  \textbf{版本号：} v1.1-dev \quad\textbar\quad \textbf{发布日期：} 2026 年 10 月 \quad\textbar\quad \textbf{文档状态：} 正式评审稿
}

% ==============================================================================
% 正文开始
% ==============================================================================
\begin{document}

\maketitle
\thispagestyle{empty}

\vspace{0.5cm}
\begin{abstract}
随着企业大规模接入多模态大语言模型（如 OpenAI、Kimi、DeepSeek 等），评估员工模型采纳率、任务结构、交互规模及端到端成功率成为内部资源调度与采购选型的重要依据；同时，部分统计结果需对外向模型供应商或产业生态进行合规反馈。然而，直接发布基于明细遥测日志的精确统计值存在通过极端活跃模式、稀疏单元交叉反推员工个体行为轨迹的重构风险。

本文提出了 \textbf{Private LLM Analytics} 系统的完整技术方案。系统确立了以单个员工所有历史记录为基础的“用户级”（User-Level）邻接关系，构建了以指标为中心的“专属差分隐私机制框架”（Statistic-Specific DP Framework），针对用户使用率（User Prevalence）、用户归一化使用占比（User-Normalized Share）、交互总量（Interaction Volume）、任务成功率（Success Rate）以及热门任务发现（Top-K Tasks）分别进行用户级贡献限制和 $L_2$ 敏感度界分析，并采用与统计结构匹配的高斯释放（Gaussian release）；其中任务成功率（Success Rate）使用成功数与总量（success / total）的单次联合高斯释放。在系统隐私记账层面，采用 Rényi 差分隐私（RDP）构建持久化、事务一致的生命周期隐私账本，提出基于语义哈希的规范发布与重放机制（Canonical Release \& Replay）以降低重复查询的噪声平均攻击风险；针对多供应商外部数据共享场景，建立了接收方作用域隔离与全局外泄合并机制；针对长周期观测需求，研发了基于时间二叉树（Binary Time Tree）的持续观测机制（Continual Observation）。实验与测试结果表明，全套 542 项工程测试套件全量通过，在 50,000 用户量级的企业级基准数据上，任务成功率百分点误差小于 0.1 pp，系统能够为企业在外部接收方环境下安全开展多模型效能度量提供具有数学严格证明的差分隐私保障。
\end{abstract}

\vspace{0.5cm}
\begin{center}
\small\textbf{关键词：} 差分隐私；大语言模型遥测；用户级敏感度；Rényi 差分隐私；持续观测；时间二叉树；数据安全合规
\end{center}

\newpage
\tableofcontents
\newpage

% 符号表与术语表
\section*{符号与术语约定 (Notation and Glossary)}
\addcontentsline{toc}{section}{符号与术语约定}

为了保证全篇技术方案在数学推导、算法描述与工程实现之间的一致性，表~\ref{tab:notation} 汇总了本文所采用的数学符号体系及对应的工程变量含义。

\begin{table}[htbp]
\centering
\small
\caption{核心数学符号与系统术语对照表}
\label{tab:notation}
\begin{tabularx}{\textwidth}{llX}
\toprule
\textbf{符号} & \textbf{工程变量 / 类型} & \textbf{数学与业务定义} \\
\midrule
$\mathcal{U}, u$ & \texttt{user\_id} (str) & 企业员工集合及单个员工个体索引 \\
$N_{\text{public}}$ & \texttt{public\_population\_size} (int) & 企业固定公开在册员工总人数（使用率计算的固定基准分母） \\
$\mathcal{P}, \mathcal{T}$ & \texttt{providers, tasks} (list) & 系统注册的模型供应商集合（基数 $P$）与业务任务集合（基数 $T$） \\
$D$ & \texttt{d\_cells} (int) & 分类法统计单元总维度，$D = P \times T$ \\
$m_u \in \mathbb{R}^D$ & \texttt{matrix[u, :]} (ndarray) & 用户 $u$ 在观测窗口内向各单元发起的原始交互调用计数向量 \\
$b_u \in \{0, 1\}^D$ & \texttt{b\_u} (ndarray) & 用户 $u$ 的二值化参与指示向量，$b_{u,j} = \mathbf{1}[m_{u,j} > 0]$ \\
$k$ & \texttt{max\_cells\_per\_user} (int) & 单用户活跃单元数 $L_0$ 范数上界（$\|b_u\|_0 \le \min(k, D)$） \\
$x_u \in \mathbb{R}^D$ & \texttt{x\_u} (ndarray) & 用户 $u$ 的归一化工作流占比向量，$x_u = m_u / \|m_u\|_1$ \\
$C$ & \texttt{clip\_norm} (float) & 交互总量统计的用户级 $L_2$ 范数贡献限制阈值 \\
$C_{\text{rate}}$ & \texttt{clip\_norm\_rate} (float) & 任务成功率联合统计中分母总量向量的 $L_2$ 截断阈值 \\
$w_u \in \mathbb{R}^{2D}$ & \texttt{w\_u} (ndarray) & 任务成功率的联合向量，$w_u = [\bar{s}_u ; \bar{n}_u]$ \\
$\tau$ & \texttt{min\_noisy\_denominator} (float) & 成功率有效发布的最小加噪分母阈值，低于该值执行抑制 \\
$\Delta_2$ & \texttt{sensitivity} (float) & 用户级加/删邻接关系下查询函数 $f$ 的全局 $L_2$ 敏感度 \\
$\sigma$ & \texttt{noise\_scale / sigma} (float) & 高斯差分隐私机制采样所采用的噪声分布标准差 \\
$\varepsilon, \delta$ & \texttt{epsilon, delta} (float) & $(\varepsilon, \delta)$-近似差分隐私的隐私损失与失效概率参数 \\
$\alpha$ & \texttt{alpha} (float / Decimal) & Rényi 差分隐私的散度阶数，满足 $\alpha \in (1, \infty)$ \\
$\rho(\alpha)$ & \texttt{rdp\_cost} (Decimal) & 阶数为 $\alpha$ 时的单次或累积 Rényi 差分隐私散度上界 \\
$\varepsilon_{\text{cap}}$ & \texttt{epsilon\_cap} (float) & 企业根据合规政策配置的全生命周期累积隐私预算硬上限 \\
$\mathcal{S}_{\text{scope}}$ & \texttt{budget\_scope\_id} (str) & 隐私记账作用域（如 \texttt{GLOBAL\_EXTERNAL}、\texttt{internal}） \\
$C_{\text{day}}, C_{\text{epoch}}$ & \texttt{day\_l2, epoch\_l12} (float) & 持续时间序列观测中的每日截断界与代际（Epoch）综合截断界 \\
$E_{\text{clip}}$ & \texttt{clipping\_bias} (float) & 贡献截断引入的结构性截断偏差，$E_{\text{clip}} = M_{\text{clip}} - M_{\text{raw}}$ \\
$E_{\text{noise}}$ & \texttt{dp\_noise\_error} (float) & 差分隐私随机噪声引入的随机扰动误差，$E_{\text{noise}} = M_{\text{DP}} - M_{\text{clip}}$ \\
$E_{\text{total}}$ & \texttt{total\_error} (float) & 最终发布值与原始真实值之间的综合总误差，$E_{\text{total}} = M_{\text{DP}} - M_{\text{raw}}$ \\
\bottomrule
\end{tabularx}
\end{table}

\paragraph{邻接关系规范描述：}
本文全文遵循标准的\textbf{用户级差分隐私（User-Level DP）}定义。若两个交互数据集 $D, D'$ 满足 $D' = D \cup \{u^*\}$ 或 $D = D' \cup \{u^*\}$（即二者差异恰好为某一个员工个体 $u^*$ 在观测窗口内的\textbf{全部交互记录集合}），则称 $D$ 与 $D'$ 互为用户级邻居数据集（Neighboring Datasets），记作 $D \sim D'$。本文所声明的所有隐私上界均在此邻接关系下严格推导并获得成立保证。


% 正文各章节
\section{项目背景与建设目标}
\label{sec:background}

\subsection{企业多模型接入与混合调用场景}
随着大语言模型（Large Language Models, LLMs）在企业知识管理、代码研发、业务分析与客服自动化等核心场景的深度渗透，现代企业的基础设施架构普遍演进为“多模型混合接入”格局。为了在推理延迟、任务能力、采购成本以及可用性冗余之间取得最优权衡，企业统一 AI 网关通常同时路由并对接多家第三方前沿模型服务商，典型代表包括专注于通用逻辑与复杂工程代码生成的 OpenAI 系列模型、擅长超长文本解析与文档精炼的 Kimi 系列模型，以及具备深度推理与专业代码能力的 DeepSeek 系列模型。

在企业实际生产环境中，员工通过统一身份认证（SSO）访问内部门户或 IDE 插件发起模型调用，系统网关实时沉淀了包含用户标识（\texttt{user\_id}）、请求时间戳（\texttt{timestamp}）、目标供应商（\texttt{provider}）、业务任务类型（\texttt{task}）以及端到端执行状态（\texttt{success}）的高频遥测日志数据（Telemetry Data）。

\subsection{企业内部运营与效能度量诉求}
为了支持企业管理层、研发效能团队、财务采购部门以及数据治理团队的精细化决策，企业迫切需要基于沉淀的遥测日志回答以下关键业务问题：
\begin{enumerate}[label=(\arabic*), leftmargin=2em]
  \item \textbf{技术采纳广度（User Prevalence）：} 全公司整体及各业务线员工对大模型能力的渗透率如何？有多少比例的员工真正将大模型引入了日常工作流？
  \item \textbf{业务结构占比（User-Normalized Share）：} 员工主要利用大模型解决哪一类工作？任务分布是以技术研发为主，还是以文档撰写与数据分析为主？
  \item \textbf{交互调用规模（Interaction Volume）：} 整体交互规模与各供应商的承载负荷是多少？各部门的算力需求呈现何种演进趋势？
  \item \textbf{服务完成质量（Success Rate）：} 不同模型供应商在特定任务类型下的端到端成功率如何？是否存在某供应商在特定业务域质量显著劣化的情况？
  \item \textbf{核心热点分布（Top-K Tasks）：} 在长周期演进中，哪些任务分类处于调用前列，值得投入自研工具链进行深度针对性优化？
\end{enumerate}

上述度量结果直接关乎算力预算分配、模型采购合同谈判权重倾斜，以及各部门 AI 辅助效能的客观评估。

\subsection{对外数据共享面临的敏感信息暴露风险}
在企业与大模型服务商的商业合作中，供应商通常要求企业反馈阶段性调用质量与业务渗透统计，以便协同调优模型、优化推理专线并制定后续支持策略。同时，行业协会、外部审计机构或科研合作方也存在对脱敏聚合统计的合法共享诉求。

然而，企业在商业协作中不能直接向外部不可信实体提供明细日志，甚至\textbf{直接发布未经隐私扰动的精确统计值同样蕴含严重的数据安全红线}：
\begin{itemize}[leftmargin=2em]
  \item \textbf{稀疏单元重构攻击（Reconstruction via Sparsity）：} 在按部门、供应商与任务交叉细分的统计表格中，部分特定业务单元（例如“某边缘任务 $\times$ 某模型”）的调用量可能仅由极少数关键员工产生。外部接收方通过分析精确数值，极易重构特定核心研发人员或小微项目组的工作轨迹。
  \item \textbf{差分推断攻击（Differencing Attack）：} 若系统持续对外发布高精度的每日或每周统计值，外部攻击者通过对相邻时间窗口的统计差值进行代数求解，能够高概率推断出某特定员工是否休假、离职或在特定敏感时段发起了密集异常调用。
  \item \textbf{商业竞争与业务架构刺探：} 精确的业务任务比例直接暴露了企业的技术研发重心、正在推进的核心业务形态，乃至企业对不同供应商调用配比的底层商业机密。
\end{itemize}

传统的去标识化方案（如移除员工姓名或对工号做单向 Hash 处理）仅阻断了浅层的身份关联，在面对拥有外部背景知识（Auxiliary Knowledge）的强攻击者时，无法抵御基于行为时序与调用特征的重识别（Re-identification）。

\subsection{系统建设目标}
针对上述业务冲突与安全合规痛点，本项目研发了 \textbf{Private LLM Analytics}（企业大模型使用数据差分隐私分析系统）。系统的核心目标概括如下：
\begin{techbox}{系统总体目标}
构建一套自主受控的企业级差分隐私统计分析基础设施，在设定的严格用户级邻接关系与 $(\varepsilon, \delta)$ 隐私参数下，为企业大模型使用日志生成具备严格数学稳定性保证的聚合统计报告；使企业既能在内部获得高保真度的业务效能度量依据，又能在无需担忧单名员工行为轨迹被外部逆向反推的前提下，安全向外部大模型供应商反馈高质量统计洞察。
\end{techbox}

\subsection{系统核心设计原则}
为了达成上述目标，系统在底层架构设计上严格遵循以下四项原则：
\begin{enumerate}[label=\textbf{原则\arabic*：}, leftmargin=4.5em]
  \item \textbf{指标专属机制（Statistic-Specific）：} 摒弃将所有业务指标简单套用单一通用敏感度与拉普拉斯/高斯加噪的做法。针对比例、占比、总量、成功率与排序等不同统计算子的取值空间，分别推导针对性的贡献限制算子、紧致敏感度与专有联合发布机制。
  \item \textbf{两层误差解耦（Two-Tier Error Decomposition）：} 在效用分析中，清晰界定“贡献截断偏差（Clipping Bias）”与“差分隐私高斯加噪误差（DP Noise Error）”，为工程参数调优与算法演进提供客观的数值归因。
  \item \textbf{规范重放防护（Canonical Replay Protection）：} 基于全维度语义哈希实现严格的幂等发布控制，相同的查询请求强制命中历史发布结果，净增隐私消耗恒等于零，避免生成独立重复噪声样本，防范攻击者通过重复查询求平均来抵消随机噪声的均值恢复攻击。
  \item \textbf{故障闭环安全（Fail-Closed Policy）：} 任何超出预设全生命周期预算上限（Lifetime Budget Cap）、违反接收方作用域策略或存在参数异常的查询，系统均无条件阻断执行，绝不向不可信接收方降级回退原始数据。
\end{enumerate}

\section{需求分析与威胁模型}
\label{sec:threat_model}

\subsection{数据来源与遥测日志特征}
系统处理的核心数据源为企业大模型网关沉淀的交互遥测记录。每条记录代表一次原子化的用户与大模型的交互会话，具备以下数据模式特征：
\begin{equation}
  r = (\texttt{interaction\_id}, \texttt{user\_id}, \texttt{timestamp}, \texttt{provider}, \texttt{task}, \texttt{success})
\end{equation}
在企业实际环境中，遥测日志呈现出显著的\textbf{用户异质性（User Heterogeneity）}与\textbf{长尾分布（Heavy-tailed Distribution）}：
\begin{itemize}[leftmargin=2em]
  \item \textbf{重度调用者（Power Users）：} 处于研发核心岗位的员工可能通过自动化脚本或频繁的代码补全插件，在单个工作日内发起数百乃至上千次调用。
  \item \textbf{轻度与偶发调用者（Casual Users）：} 部分职能部门员工可能在整月周期内仅发起个位数的长文本摘要或咨询交互。
  \item \textbf{专长偏好分化（Persona Bias）：} 员工在任务类型与模型供应商的选择上具备高度稳定的个人习惯（例如某些工程师固定偏好 DeepSeek 解决算法题目，而文档专员固定使用 Kimi 处理会议纪要）。
\end{itemize}

\subsection{统计对象与多维空间划分}
系统将统计分析的目标域形式化为二维笛卡尔积空间：
\begin{equation}
  \mathcal{C} = \mathcal{P} \times \mathcal{T} = \{ (\text{prov}_p, \text{task}_t) \mid p \in \{1, \dots, P\}, \, t \in \{1, \dots, T\} \}
\end{equation}
统计单元总数记为 $D = |\mathcal{C}| = P \times T$。所有发布的统计指标均在该 $D$ 维离散单元网格上定义。

\subsection{隐私保护对象与保护粒度}
在隐私保护的粒度选择上，差分隐私存在“记录级”（Record-Level）与“用户级”（User-Level）两种不同的邻接关系定义：
\begin{itemize}[leftmargin=2em]
  \item \textbf{记录级保护对高频用户的保障局限：} 记录级邻接关系仅约束相邻数据集相差“单条调用记录”。根据群隐私性质（Group Privacy）~\cite{dwork2014algorithmic}，当将记录级差分隐私扩展至一个在观测周期内包含 $k$ 条交互记录的员工个体时，其用户级保护界会随记录数 $k$ 显著恶化（隐私损失参数放大约至 $k\varepsilon$；而在近似差分隐私下，$\delta$ 项也会放大至约 $k e^{(k-1)\varepsilon} \delta$）。对于拥有高频交互特征的企业员工，记录级差分隐私无法提供具有实用意义的用户级保护界。
  \item \textbf{用户级邻接关系的直接采用：} 本系统直接确立\textbf{用户级邻接关系（User-Level Adjacency）}，将隐私保护的原子单位定义为单个自然人员工在整个观测窗口内的完整行为轨迹，确保隐私保证与单用户的交互调用次数解耦。
\end{itemize}

\begin{definitionbox}{用户级差分隐私 (User-Level Differential Privacy)}
设 $\mathcal{M}$ 为一个随机化发布算法。若对于任意两个在观测时间窗口内相差某个特定员工 $u^*$ 全部调用记录的邻近数据集 $D, D'$（即 $D \sim D'$），以及算法输出空间中的任意可测子集 $\mathcal{S}$，均满足：
\begin{equation}
  \Pr[\mathcal{M}(D) \in \mathcal{S}] \le e^\varepsilon \Pr[\mathcal{M}(D') \in \mathcal{S}] + \delta
\end{equation}
则称算法 $\mathcal{M}$ 满足用户级 $(\varepsilon, \delta)$-差分隐私。
\end{definitionbox}

该数学保证向员工提供了明确的隐私底线：无论外部实体具备何种分析能力，观测到发布结果的概率分布对于“该员工是否使用大模型”或“该员工全部使用轨迹的存在与否”保持极高程度的平滑稳定性。

\subsection{外部接收方威胁模型与能力假设}
系统面向的潜在威胁实体包括外部模型供应商、第三方合作机构，以及通过公开或半公开渠道截获统计报表的潜在恶意观测者。系统对威胁模型做出如下界定：
\begin{enumerate}[label=(\alph*), leftmargin=2.5em]
  \item \textbf{半诚实与恶意外部接收方：} 接收方在遵守基本协议规范的同时，会利用自身算力与数学工具，极尽所能从接收到的统计报告中挖掘企业核心成员的行为特征与组织架构信息。
  \item \textbf{任意外部背景知识（Arbitrary Auxiliary Information）：} 依据差分隐私的严密理论基础~\cite{dwork2006calibrating}，假设攻击者掌握除目标被保护数据之外的一切辅助信息（例如员工在社交平台公开的代码提交记录、离职变动公示、外部公开的技术博客等）。差分隐私在此类强假设下依然保持稳健。
  \item \textbf{任意后处理能力（Post-Processing Immunity）：} 接收方可对发布结果应用任何确定性或随机化的后续统计变换、机器学习建模或数值反演算法。根据差分隐私后处理定理，此类操作绝不可能凭空削弱已成立的隐私保证。
\end{enumerate}

\subsection{外部合谋假设与作用域隔离}
在多供应商共存的商业生态中，不能简单假定各供应商之间彼此信息绝缘。系统引入了审慎的\textbf{全局外泄合谋假设（Collusion Assumption）}：
\begin{techbox}{全局外部合谋假设 (Global External Collusion Assumption)}
系统默认假设所有外部模型供应商（如 OpenAI、Kimi、DeepSeek 等）可能在离线环境下共享其收到的企业统计报表，或被同一外部第三方咨询/评测机构集中获取。因此，系统将所有属于不可信外部范畴的接收方统一绑定至同一个全局隐私预算作用域（\texttt{GLOBAL\_EXTERNAL}）。任何外部接收方发起的新发布均累积计入该共享作用域的终身消耗。
\end{techbox}

对于企业内部经过严格鉴权的高信任治理团队（\texttt{internal}），由于其肩负算法效能审计与真实效用评测职责，系统为其开辟独立的内部预算作用域，与外部暴露严格解耦。

\subsection{系统信任边界 (Trust Boundary)}
系统采用经典的\textbf{可信数据保管者模型（Trusted Curator Model）}。系统的信任边界划定如图~\ref{fig:trust_boundary} 所示：
\begin{itemize}[leftmargin=2em]
  \item \textbf{可信边界内部（Trusted Computing Boundary）：} 企业私有部署的遥测日志数据库、数据预聚合引擎、贡献限制算子、噪声采样器、RDP 记账账本与状态持久化引擎。该区域受企业信任边界保护，实施严格的网络隔离与访问控制。
  \item \textbf{不可信外部边界（Untrusted Consumers）：} 外部模型供应商 API 回传接口、公开 CSV/JSON 报表下载端点以及面向非特权租户的展示界面。此区域仅允许流经差分隐私扰动后的数值产品。
\end{itemize}

\begin{figure}[htbp]
\centering
\begin{tikzpicture}[
  scale=0.9, every node/.style={scale=0.9},
  box/.style={draw=navyblue, fill=lightbg, rounded corners=3pt, thick, align=center, minimum height=1cm, inner sep=6pt},
  trust/.style={draw=accentgreen, fill=accentgreen!5, rounded corners=6pt, dashed, thick},
  untrust/.style={draw=accentamber, fill=accentamber!5, rounded corners=6pt, dashed, thick},
  arrow/.style={-{Latex[length=3mm]}, thick, color=navyblue}
]

  % 可信区
  \node[box, text width=3cm] (raw) at (0, 0) {企业明细遥测日志\\(\texttt{Raw Telemetry})};
  \node[box, text width=3.2cm] (curator) at (4.5, 0) {可信隐私计算核心\\(Clipping + DP Noise)};
  \node[box, text width=3cm] (ledger) at (4.5, -2.5) {全生命周期账本\\(\texttt{RDP Ledger})};

  % 边界外
  \node[box, text width=3.2cm] (release) at (9.5, 0) {发布与导出接口\\(DP Estimates)};
  \node[box, text width=2.8cm, fill=white] (extA) at (14, 1.2) {OpenAI\\(不可信接收方)};
  \node[box, text width=2.8cm, fill=white] (extB) at (14, 0) {Kimi\\(不可信接收方)};
  \node[box, text width=2.8cm, fill=white] (extC) at (14, -1.2) {DeepSeek\\(不可信接收方)};

  % 连线
  \draw[arrow] (raw) -- (curator);
  \draw[arrow] (curator) -- (ledger);
  \draw[arrow] (curator) -- node[above]{\small 仅限DP发布值} (release);
  \draw[arrow] (release) -- (extA);
  \draw[arrow] (release) -- (extB);
  \draw[arrow] (release) -- (extC);

  % 信任区域框
  \begin{scope}[on background layer]
    \node[trust, fit=(raw) (curator) (ledger), label={[accentgreen, font=\bfseries]above:企业可信安全计算边界 (Trusted Curator Boundary)}] {};
    \node[untrust, fit=(release) (extA) (extB) (extC), label={[accentamber, font=\bfseries]above:不可信外部暴露域 (Untrusted External Domain)}] {};
  \end{scope}

\end{tikzpicture}
\caption{系统信任边界与数据流转拓扑示意图}
\label{fig:trust_boundary}
\end{figure}

\subsection{系统非保护目标与技术边界澄清}
为了在学术、合规及架构评审中确立严谨的技术预期，明确系统\textbf{不保护}的内容：
\begin{enumerate}[label=(\arabic*), leftmargin=2em]
  \item \textbf{宏观总体统计事实（Macro Population Properties）：} 差分隐私旨在掩盖个体存在性与微观贡献，但\textbf{其数学设计的本义就是要忠实反映宏观群体规律}。例如，“全公司整体有 90\% 的员工积极使用了大模型”属于宏观统计属性，系统将准确发布这一事实，该宏观事实不属于差分隐私保护的个体隐私范畴。
  \item \textbf{替代身份认证与网络访问控制：} 差分隐私是针对统计发布与数据分析的数学扰动机制，绝非替代网络防火墙、传输层加密（TLS）或细粒度角色访问控制（RBAC）的工具。
  \item \textbf{抗侧信道硬件泄露：} 假定底层计算硬件与操作系统满足标准的内存与计算隔离安全假设，非硬件侧信道攻击防护方案。
\end{enumerate}

\section{系统总体架构}
\label{sec:architecture}

\subsection{分层架构设计与核心组件}
为了实现差分隐私数学机制、全生命周期记账与高并发企业级服务之间的解耦，\textbf{Private LLM Analytics} 采用了模块化、高内聚的分层架构体系。整个系统自底向上划分为六个核心层级，如图~\ref{fig:architecture} 所示。

\begin{figure}[htbp]
\centering
\begin{tikzpicture}[
  scale=0.88, every node/.style={scale=0.88},
  layerbox/.style={draw=navyblue, fill=lightbg, rounded corners=4pt, thick, inner sep=8pt, text width=14cm},
  subbox/.style={draw=slateblue, fill=white, rounded corners=2pt, font=\small, align=center, inner sep=4pt, minimum height=0.8cm},
  arrow/.style={-{Latex[length=2.5mm]}, thick, color=navyblue}
]

  % Layer 1: Ingestion
  \node[layerbox] (L1) at (0, 0) {
    \textbf{\color{navyblue} 1. 遥测日志接入与校验层 (Data Ingestion \& Validation Layer)}\\[3pt]
    \begin{tikzpicture}
      \node[subbox, text width=3.8cm] at (0,0) {CSV / Parquet 校验解析\\(\texttt{Schema Validation})};
      \node[subbox, text width=3.8cm] at (4.4,0) {流式 SHA-256 数据指纹\\(\texttt{Dataset Fingerprint})};
      \node[subbox, text width=3.8cm] at (8.8,0) {动态分类学元数据映射\\(\texttt{Dynamic Taxonomy})};
    \end{tikzpicture}
  };

  % Layer 2: Aggregation
  \node[layerbox, below=0.45cm of L1] (L2) {
    \textbf{\color{navyblue} 2. 用户级状态矩阵聚合层 (User Matrix Aggregation Layer)}\\[3pt]
    \begin{tikzpicture}
      \node[subbox, text width=3.8cm] at (0,0) {单用户调用量聚合矩阵\\$\mathbf{M} \in \mathbb{R}^{N \times D}$};
      \node[subbox, text width=3.8cm] at (4.4,0) {单用户成功量聚合矩阵\\$\mathbf{S} \in \mathbb{R}^{N \times D}$};
      \node[subbox, text width=3.8cm] at (8.8,0) {固定公开人口基准绑定\\$N_{\text{public}}$};
    \end{tikzpicture}
  };

  % Layer 3: Statistic Profile
  \node[layerbox, below=0.45cm of L2] (L3) {
    \textbf{\color{navyblue} 3. 统计专属机制层 (Statistic-Specific DP Framework Layer)}\\[3pt]
    \begin{tikzpicture}
      \node[subbox, text width=2.4cm] at (0,0) {User Prevalence\\($L_0$ 截断, $\Delta_2 \le \sqrt{k}$)};
      \node[subbox, text width=2.5cm] at (2.8,0) {Normalized Share\\(单纯形映射, $\Delta_2 \le 1$)};
      \node[subbox, text width=2.5cm] at (5.7,0) {Interaction Volume\\($L_2$ 截断, $\Delta_2 = C$)};
      \node[subbox, text width=2.5cm] at (8.6,0) {Success Rate\\(联合向量, $\Delta_2 \le \sqrt{2}C_{\text{rate}}$)};
      \node[subbox, text width=2.1cm] at (11.3,0) {Top-K Tasks\\(全字典加噪)};
    \end{tikzpicture}
  };

  % Layer 4: Mechanism & Sampling
  \node[layerbox, below=0.45cm of L3] (L4) {
    \textbf{\color{navyblue} 4. 差分隐私机制与参数校准层 (Mechanism \& Calibration Layer)}\\[3pt]
    \begin{tikzpicture}
      \node[subbox, text width=3.8cm] at (0,0) {解析高斯校准引擎\\(\texttt{Analytic Gaussian})};
      \node[subbox, text width=3.8cm] at (4.4,0) {操作系统物理真随机数控制\\(\texttt{OS Entropy RNG})};
      \node[subbox, text width=3.8cm] at (8.8,0) {持续观测时间二叉树机制\\(\texttt{Binary Time Tree})};
    \end{tikzpicture}
  };

  % Layer 5: Accounting & Governance
  \node[layerbox, below=0.45cm of L4] (L5) {
    \textbf{\color{navyblue} 5. 全生命周期记账与发布控制层 (Accounting \& Governance Layer)}\\[3pt]
    \begin{tikzpicture}
      \node[subbox, text width=3.8cm] at (0,0) {Rényi 差分隐私紧致组合\\(\texttt{RDP Accountant})};
      \node[subbox, text width=3.8cm] at (4.4,0) {规范重放零增量缓存引擎\\(\texttt{Canonical Replay})};
      \node[subbox, text width=3.8cm] at (8.8,0) {接收方作用域隔离与账本持久化\\(\texttt{SQLite State Manager})};
    \end{tikzpicture}
  };

  % Layer 6: Application & Presentation
  \node[layerbox, below=0.45cm of L5] (L6) {
    \textbf{\color{navyblue} 6. 业务应用与可解释性导出层 (Application \& Presentation Layer)}\\[3pt]
    \begin{tikzpicture}
      \node[subbox, text width=3.8cm] at (0,0) {合规发布脱敏 CSV/JSON 导出\\(\texttt{Secure Exporter})};
      \node[subbox, text width=3.8cm] at (4.4,0) {差分隐私 Web 控制台向导\\(\texttt{Interactive Console})};
      \node[subbox, text width=3.8cm] at (8.8,0) {内部效用对比与 $\varepsilon$-Sweep 评估\\(\texttt{Utility Evaluator})};
    \end{tikzpicture}
  };

  % 连接线
  \draw[arrow] (L1) -- (L2);
  \draw[arrow] (L2) -- (L3);
  \draw[arrow] (L3) -- (L4);
  \draw[arrow] (L4) -- (L5);
  \draw[arrow] (L5) -- (L6);

\end{tikzpicture}
\caption{Private LLM Analytics 系统六层技术架构图}
\label{fig:architecture}
\end{figure}

\subsection{核心层级职能描述}
各层级具体职能与设计考量说明如下：
\begin{enumerate}[label=\textbf{\arabic*.}, leftmargin=2em]
  \item \textbf{遥测日志接入与校验层：} 支持百万级规模的 CSV 与 Parquet 格式交互遥测数据输入。负责严格检验必需列的存在性与类型合法性；采用文件流与内容双重 SHA-256 哈希计算数据集指纹，确保所有后续发布均与明确的数据集快照版本形成强绑定。
  \item \textbf{用户级状态矩阵聚合层：} 依托高度优化的矢量化算子，将离散日志快速投影为以单个员工为维度的稀疏/密集矩阵 \texttt{UserMatrix}。该层负责将每个用户的全部离散行为聚合为状态向量，奠定用户级差分隐私的数学输入基础。
  \item \textbf{统计专属机制层：} 系统的核心创新组件。针对不同业务查询，实例化专有的统计 Profile（如 \texttt{UserPrevalenceProfile}、\texttt{SuccessRateProfile} 等）。由各 Profile 自主负责实现该统计指标的用户级贡献限制（Contribution Bounding）和理论敏感度求取。
  \item \textbf{差分隐私机制与参数校准层：} 基于精确的解析高斯机制（Analytic Gaussian Mechanism）求解最优噪声标准差 $\sigma$；生产模式下严密阻断任何调用方对随机数发生器（RNG）的控制，确保噪声样本源自操作系统提供的熵源（OS-provided entropy）初始化的随机数生成器。对于流式时间序列，该层挂载时间二叉树机制实现增量加噪。
  \item \textbf{全生命周期记账与发布控制层：} 维护持久化、事务一致的生命周期隐私账本；正常服务重启不会自动重置历史状态。通过对比查询的规范语义哈希判断是否存在历史结果：命中则直接返回已有结果且增量预算为零（规范重放）；未命中则执行原子预留（Compare-and-Reserve）与多阶 Rényi 差分隐私累加记账。
  \item \textbf{业务应用与可解释性导出层：} 负责对加噪后的向量执行合规后处理（如值域截断、归一化与阈值抑制）；向外部接收方提供纯脱敏的 CSV 与 JSON 结构化数据下载；向企业内部管理员提供两层误差拆解与 $\varepsilon$ 效用敏感度扫描报告。
\end{enumerate}

\subsection{端到端请求流转与状态时序}
一次典型的差分隐私发布请求在系统内部经历如图~\ref{fig:flowchart} 所示的状态生命周期。

\begin{figure}[htbp]
\centering
\begin{tikzpicture}[
  scale=0.88, every node/.style={scale=0.88},
  node distance=1.2cm,
  startstop/.style={draw=navyblue, fill=navyblue!10, rounded corners=3pt, thick, align=center, inner sep=5pt, text width=3.6cm},
  process/.style={draw=slateblue, fill=white, rounded corners=2pt, thick, align=center, inner sep=5pt, text width=3.8cm},
  decision/.style={draw=accentamber, fill=accentamber!10, diamond, aspect=2, thick, align=center, inner sep=2pt, text width=2.4cm},
  arrow/.style={-{Latex[length=2.5mm]}, thick, color=navyblue}
]

  \node[startstop] (start) {收到查询请求\\(\texttt{ReleaseRequest})};
  \node[process, below=0.8cm of start] (hash) {计算规范语义哈希\\(\texttt{Canonical Key})};
  \node[decision, below=0.8cm of hash] (cache) {存在完全匹配\\的历史发布？};
  
  % 缓存分支
  \node[process, right=2.2cm of cache] (replay) {规范重放 (Canonical Replay)\\$\Delta\text{RDP} = 0$, 零预算消耗};
  \node[startstop, below=3.5cm of replay] (finish) {返回发布结果与导出视图\\(\texttt{DP Estimate})};

  % 未命中分支
  \node[process, below=1.0cm of cache] (reserve) {原子预算预留校验\\$\varepsilon_{\text{current}} + \Delta\varepsilon \le \varepsilon_{\text{cap}}$};
  \node[decision, below=0.8cm of reserve] (budget) {预算校验\\通过？};
  \node[process, left=2.2cm of budget] (fail) {\color{red}\bfseries 拒绝服务\\(\texttt{BudgetExceeded})\\Fail-Closed 阻断};
  
  \node[process, below=1.0cm of budget] (exec) {执行统计专属机制\\(截断 + 解析高斯加噪)};
  \node[process, below=0.8cm of exec] (commit) {提交 RDP 账本与生成缓存\\(\texttt{Commit Release})};

  % 连线
  \draw[arrow] (start) -- (hash);
  \draw[arrow] (hash) -- (cache);
  \draw[arrow] (cache) -- node[above]{\small 命中} (replay);
  \draw[arrow] (replay) -- (finish);
  
  \draw[arrow] (cache) -- node[right]{\small 未命中} (reserve);
  \draw[arrow] (reserve) -- (budget);
  \draw[arrow] (budget) -- node[above]{\small 否} (fail);
  \draw[arrow] (budget) -- node[right]{\small 是} (exec);
  \draw[arrow] (exec) -- (commit);
  \draw[arrow] (commit) |- (finish);

\end{tikzpicture}
\caption{差分隐私查询原子生命周期与安全流转时序图}
\label{fig:flowchart}
\end{figure}

\section{数据模型与统计任务}
\label{sec:data_statistics}

\subsection{输入数据模式与字段语义}
系统对接入的企业大模型交互遥测日志执行强类型模式（Schema）校验。必需字段及其业务语义与约束条件如表~\ref{tab:input_schema} 所示。

\begin{table}[htbp]
\centering
\small
\caption{企业大模型交互遥测日志模式规范}
\label{tab:input_schema}
\begin{tabularx}{\textwidth}{llcX}
\toprule
\textbf{字段名称} & \textbf{数据类型} & \textbf{可空性} & \textbf{约束条件与业务语义说明} \\
\midrule
\texttt{interaction\_id} & String (UUID) & 否 & 单次大模型调用的全局唯一标识符，用于去重审计 \\
\texttt{user\_id} & String & 否 & 员工唯一身份标识（如企业统一工号或 SSO 标识），用于用户级分组 \\
\texttt{timestamp} & Timestamp & 否 & 交互发生时的 UTC 时间戳（ISO 8601），用于时间窗口划分 \\
\texttt{provider} & String & 否 & 调用的目标模型供应商内部标识（如 \texttt{provider\_a}、\texttt{provider\_b}） \\
\texttt{task} & String & 否 & 业务任务类别标识（如 \texttt{technical}、\texttt{document}、\texttt{analysis}） \\
\texttt{success} & Boolean & 是 & 任务端到端完成状态（True 为成功，False 为失败），成功率统计必需 \\
\bottomrule
\end{tabularx}
\end{table}

\subsection{动态分类学体系 (Dynamic Taxonomy)}
传统的差分隐私分析代码往往将统计维度硬编码为固定大小的矩阵（如历史版本遗留的 $3 \times 3 = 9$ 单元）。本项目在架构层面重构为完全参数化的\textbf{动态分类学管理机制（Dynamic Taxonomy）}。
系统允许在不修改底层差分隐私算子实现的前提下，通过修改配置文件动态增删模型供应商与任务分类。设注册的供应商集合为 $\mathcal{P} = \{p_1, \dots, p_P\}$，任务集合为 $\mathcal{T} = \{t_1, \dots, t_T\}$，分类系统自动构建双射映射：
\begin{equation}
  \Phi: \mathcal{P} \times \mathcal{T} \to \{0, 1, \dots, D-1\}, \quad \Phi(p_i, t_j) = i \cdot T + j
\end{equation}
其中全局统计维度 $D = P \times T$ 保持动态计算。系统内所有的截断算子、敏感度计算、高斯协方差矩阵以及后处理均匀回退向量均严格依赖动态传入的维度 $D$，有效防止了硬编码常数带来的维度越界与统计口径失配。

\subsection{供应商内部标识与展示名称的解耦}
为了兼顾外部合规审计的稳定性与 Web 界面的人性化展示，系统在工程上实施了\textbf{底层标识与展示视图的严格分离}：
\begin{itemize}[leftmargin=2em]
  \item \textbf{底层规范标识（Canonical ID）：} 数据库、RDP 记账账本、语义哈希计算以及接收方鉴权逻辑，均严格使用不可变的规范标识符（如 \texttt{provider\_a}、\texttt{provider\_b}、\texttt{provider\_c}、\texttt{internal}）。
  \item \textbf{展示层友好映射（Display Name）：} 仅在向 Web 控制台以及面向业务人员的最终报表输出时，执行固定的映射投影：
    $$\texttt{provider\_a} \mapsto \text{OpenAI}, \quad \texttt{provider\_b} \mapsto \text{Kimi}, \quad \texttt{provider\_c} \mapsto \text{DeepSeek}$$
\end{itemize}
这一设计确保了管理人员未来调整供应商对外展示名称时，\textbf{绝不会破坏历史 RDP 累积账本的一致性，也不会引起规范发布哈希漂移}。

\subsection{用户级交互矩阵 (\texttt{UserMatrix}) 建模}
在数据摄取阶段，系统通过高效的多列分组与稠密矩阵投影，将原始日志抽象为数据结构 \texttt{UserMatrix}。设当前活跃用户集合为 $\mathcal{U}_{\text{active}}$（大小为 $n_{\text{active}}$），企业在册固定公开总人口为 $N_{\text{public}}$，则 \texttt{UserMatrix} 包含：
\begin{equation}
  \mathbf{M} \in \mathbb{R}^{n_{\text{active}} \times D}, \quad \mathbf{M}_{u, j} = \sum_{r \in \mathcal{R}_u} \mathbf{1}[\Phi(\text{prov}_r, \text{task}_r) = j]
\end{equation}
其中 $\mathcal{R}_u$ 表示用户 $u$ 在选定时间窗口内的所有交互记录集合。若遥测日志包含 \texttt{success} 字段，系统同步构建成功量矩阵 $\mathbf{S} \in \mathbb{R}^{n_{\text{active}} \times D}$：
\begin{equation}
  \mathbf{S}_{u, j} = \sum_{r \in \mathcal{R}_u} \mathbf{1}[\Phi(\text{prov}_r, \text{task}_r) = j \land \text{succ}_r = \text{True}]
\end{equation}
在构造阶段，系统强制验证约束 $n_{\text{active}} \le N_{\text{public}}$；若活跃用户数超过公开总人口，系统执行 Fail-Closed 熔断报错，防范因配置不匹配引发的分母越界。

\subsection{五类统计任务总览导航表}
系统支持的五类统计分析任务、对应的数学贡献形式、截断机制、敏感度与效用评估度量汇总于表~\ref{tab:statistic_master}。该表作为全篇差分隐私机制设计的核心基准。

\begin{table}[htbp]
\centering
\small
\caption{Private LLM Analytics 五大核心统计任务设计总览表}
\label{tab:statistic_master}
\begin{tabularx}{\textwidth}{p{2.2cm}Xp{2.6cm}p{2.0cm}p{2.2cm}p{2.1cm}}
\toprule
\textbf{统计指标} & \textbf{业务度量目标} & \textbf{用户级贡献定义} & \textbf{截断方式} & \textbf{$L_2$ 敏感度} & \textbf{主效用度量} \\
\midrule
\textbf{用户使用率}\newline (Prevalence) & 统计在册员工中至少使用一次特定模型任务能力的比例 & 二值指示向量\newline $b_{u,j} = \mathbf{1}[m_{u,j} > 0]$ & $L_0$ 范数截断\newline ($\|b_u\|_0 \le k$) & $\sqrt{\min(k, D)}$\newline (实现为 $\sqrt{k}$) & 百分点误差 (pp)\newline $E = |\hat{r} - r_{\text{raw}}|$ \\
\addlinespace
\textbf{使用结构占比}\newline (Normalized Share) & 统计宏观任务分布，反映员工在各能力域的工作流时间分配 & 归一化工作流向量\newline $x_u = m_u / \|m_u\|_1$ & 单纯形天然有界\newline ($\|x_u\|_2 \le 1$) & $\Delta_2 \le 1.0$ & 曼哈顿距离 $L_1$\newline 及百分点误差 \\
\addlinespace
\textbf{交互总量}\newline (Volume) & 统计各模型任务的累计交互调用总规模，指导算力与成本核算 & 多维计数向量\newline $m_u \in \mathbb{R}^D$ & $L_2$ 范数截断\newline (阈值 $C$) & $\Delta_2 = C$ & 平均相对误差\newline (MRE) \\
\addlinespace
\textbf{任务成功率}\newline (Success Rate) & 评估特定模型在特定业务域的端到端调用质量与可用性 & 联合向量\newline $w_u = [\bar{s}_u ; \bar{n}_u] \in \mathbb{R}^{2D}$ & 共享比例截断\newline (总量阈值 $C_{\text{rate}}$) & $\Delta_2 \le \sqrt{2}C_{\text{rate}}$ & 百分点误差 (pp)\newline 小分母安全抑制 \\
\addlinespace
\textbf{热门任务}\newline (Top-K Tasks) & 发现企业内部调用最频繁的热门业务分类前 $K$ 名 & 任务二值指示\newline $v_{u,t} = \mathbf{1}[m_{u,t} > 0]$ & 任务数截断\newline ($\|v_u\|_0 \le k_{\text{task}}$) & $\sqrt{\min(k_t, T)}$\newline (实现为 $\sqrt{k_t}$) & Top-$K$ 命中率\newline 与 Jaccard 指数 \\
\bottomrule
\end{tabularx}
\end{table}

\section[统计专属差分隐私机制框架]{统计专属差分隐私机制框架\\[0.2cm](Statistic-Specific DP Framework)}
\label{sec:dp_profiles}

\subsection{机制设计动机：为何否定通用单一机制}
在传统差分隐私工程实践中，一种常见的简化方式是将所有统计查询统一视为数值型向量查询，不加区分地对各类统计值施加通用的全局敏感度界限与单一拉普拉斯或高斯噪声。这一方案在真实企业业务场景下暴露出严重的缺陷：
\begin{enumerate}[label=(\arabic*), leftmargin=2em]
  \item \textbf{值域约束与敏感度错配：} 比例指标（Prevalence, Success Rate）天然位于 $[0, 1]$ 区间，而调用总量（Volume）可能高达数十万次。若使用统一机制，会导致比例指标被过大噪声彻底湮没，或导致总量指标的敏感度被严重低估。
  \item \textbf{比值型统计的敏感性与不稳定性：} 成功率是成功次数与调用总次数的比值（Ratio）。直接经验比率在小分母情况下具有较高单用户敏感性和严重效用不稳定性；若将分子分母独立加噪且独立记账，不仅会消耗双倍隐私预算，而且在除法后处理时极易产生严重失真的极端野值。因此，系统联合发布受限的 success / total counts，再执行安全阈值抑制和比率后处理。
\end{enumerate}

因此，本项目提出了\textbf{统计专属差分隐私机制框架（Statistic-Specific DP Framework）}。每个注册的统计指标拥有完整的闭环 Profile：
\begin{center}
  \textbf{业务统计目标} $\longrightarrow$ \textbf{用户级贡献定义} $\longrightarrow$ \textbf{贡献截断算子} $\longrightarrow$ \textbf{紧致敏感度分析} $\longrightarrow$ \textbf{专属加噪释放} $\longrightarrow$ \textbf{合规后处理}
\end{center}

\subsection{用户使用率机制 (User Prevalence Profile)}
\paragraph{1. 业务目标与贡献定义：}
度量在册企业员工对特定模型任务单元的能力采纳渗透率。对于用户 $u$ 和统计单元 $j$，其原始贡献为二值化指示函数：
\begin{equation}
  b_{u, j} = \mathbf{1}[m_{u, j} > 0] \in \{0, 1\}
\end{equation}
无论用户 $u$ 在单元 $j$ 调用了 1 次还是 10,000 次，其对单元 $j$ 的贡献严格恒等为 1。

\paragraph{2. 贡献限制算子 ($L_0$ Bounding)：}
为了防止跨领域重度活跃员工在过多统计单元中同时产生贡献，系统施加 $L_0$ 范数上界约束 $\|b_u\|_0 \le k$。若某个用户的活跃单元数超出了 $k$，系统采用确定性用户级伪随机抽样算子，无放回截断保留 $k$ 个单元，其余置零。

\paragraph{3. 敏感度推导与紧致性：}
设全局聚合直方图为 $H(D) = \sum_{u \in \mathcal{U}} b_u$。考虑相邻数据集 $D \sim D'$（相差用户 $u^*$）：
\begin{equation}
  \|H(D) - H(D')\|_2 = \|b_{u^*}\|_2 = \sqrt{\sum_{j=1}^D b_{u^*, j}^2} = \sqrt{\sum_{j=1}^D b_{u^*, j}} = \sqrt{\|b_{u^*}\|_0} \le \sqrt{\min(k, D)}
\end{equation}
\begin{techbox}{Prevalence 敏感度定理}
在 $L_0$ 贡献限制参数为 $k$、分类单元总数为 $D$ 的前提下，用户级加/删邻接关系下直方图聚合的全局 $L_2$ 敏感度紧致上界为：
\begin{equation}
  \Delta_2^{\text{tight}} = \sqrt{\min(k, D)}
\end{equation}
在工程生产配置中，若分类维度满足 $D \ge k$，代码采用保守校准敏感度 $\Delta_2^{\text{prod}} = \sqrt{k}$（例如在默认 $k=D=9$ 的场景下，$\Delta_2 = 3.0$）。
\end{techbox}

\paragraph{4. 高斯加噪与后处理：}
系统向聚合直方图注入独立高斯噪声向量：$\tilde{H} = H(D) + \mathcal{N}(0, \sigma^2 \mathbf{I}_D)$。后处理算子严格遵循：
\begin{equation}
  \hat{H}_j = \max(0, \tilde{H}_j), \qquad \hat{r}_j = \operatorname{clip}\left(\frac{\hat{H}_j}{N_{\text{public}}}, 0.0, 1.0\right)
\end{equation}
其中 $N_{\text{public}}$ 严格采用企业固定公开在册人口，严禁使用加噪后的或动态估计的活跃人数作为分母，防范二次信息泄露。

\subsection{用户归一化占比机制 (User-Normalized Share Profile)}
\paragraph{1. 业务目标与贡献定义：}
度量宏观工作流结构分布，反映企业员工在不同模型与任务之间的精力分配倾向。为了消除“单个狂热用户主导全局占比”的偏差，系统首先在单用户内部进行自身归一化：
\begin{equation}
  x_{u, j} = \begin{cases}
    \dfrac{m_{u, j}}{\|m_u\|_1} = \dfrac{m_{u, j}}{\sum_{i=1}^D m_{u, i}}, & \text{若 } \|m_u\|_1 > 0 \\[8pt]
    0, & \text{若 } \|m_u\|_1 = 0
  \end{cases}
\end{equation}
该映射将每个活跃员工的交互偏好投影至标准单纯形 $\Delta^{D-1}$ 上。

\paragraph{2. 敏感度推导与单纯形有界性：}
对于任意活跃员工 $u$，由于 $x_{u, j} \ge 0$ 且 $\sum_{j=1}^D x_{u, j} = 1$，其 $L_1$ 范数恒满足 $\|x_u\|_1 = 1$。由向量范数关系式 $\|x_u\|_2 \le \|x_u\|_1$：
\begin{equation}
  \|x_u\|_2 = \sqrt{\sum_{j=1}^D x_{u, j}^2} \le \sum_{j=1}^D x_{u, j} = 1.0
\end{equation}
因此，对于聚合向量 $X(D) = \sum_{u \in \mathcal{U}} x_u$，增加或移除一个用户产生的最大 $L_2$ 范数偏移为：
\begin{equation}
  \Delta_2 = \max_{u^*} \|x_{u^*}\|_2 \le 1.0
\end{equation}
该指标天然具备严格的 $L_2$ 敏感度上界 $\Delta_2 \le 1.0$，无需任何人为截断算子介入。

\paragraph{3. 后处理（非负截断与单纯形重新投影）：}
对加噪后的向量 $\tilde{X} = X(D) + \mathcal{N}(0, \sigma^2 \mathbf{I}_D)$，后处理依次执行非负截断与再归一化：
\begin{equation}
  \bar{s}_j = \max(0, \tilde{X}_j), \qquad \hat{s}_j = \begin{cases}
    \dfrac{\bar{s}_j}{\sum_{i=1}^D \bar{s}_i}, & \text{若 } \sum_{i=1}^D \bar{s}_i > 0 \\[8pt]
    \dfrac{1}{D}, & \text{若 } \sum_{i=1}^D \bar{s}_i = 0
  \end{cases}
\end{equation}
此方法保证了发布向量严格满足 $\sum_{j=1}^D \hat{s}_j = 1$ 且 $\hat{s}_j \ge 0$。

\subsection{交互总量机制 (Interaction Volume Profile)}
\paragraph{1. 业务目标与贡献限制：}
度量各业务单元的绝对交互规模，直接指导基础设施算力采购与云账单核算。由于用户的原始调用计数 $m_u$ 理论上无自然上界，系统采用标准的用户级 $L_2$ 截断算子（$L_2$ Clipping）：
\begin{equation}
  \bar{m}_u = m_u \cdot \min\left(1, \frac{C}{\|m_u\|_2}\right)
\end{equation}
其中 $C > 0$ 为配置的用户级最大贡献阈值。

\paragraph{2. 敏感度与两层误差分解：}
截断后的单用户向量必然满足 $\|\bar{m}_u\|_2 \le C$，因此在用户级邻接关系下的全局 $L_2$ 敏感度为：
\begin{equation}
  \Delta_2 = C
\end{equation}
向截断后的聚合值 $\bar{M} = \sum_{u \in \mathcal{U}} \bar{m}_u$ 注入高斯噪声：$\hat{M} = \bar{M} + \mathcal{N}(0, \sigma^2 \mathbf{I}_D)$。发布值 $\hat{M}$ 与未受保护的原始真实值 $M = \sum_{u \in \mathcal{U}} m_u$ 之间的总误差具备精确的代数分解结构：
\begin{equation}
  \label{eq:volume_error_decomp}
  E_{\text{total}} = \hat{M} - M = \underbrace{(\hat{M} - \bar{M})}_{E_{\text{noise}} \sim \mathcal{N}(0, \sigma^2)} + \underbrace{(\bar{M} - M)}_{E_{\text{clip}} \le 0}
\end{equation}
\begin{techbox}{误差解耦洞察}
在企业高倾斜的长尾分布下，重度员工发起的调用量往往远超预设的常数阈值（如 $C=20$）。系统量化评估证明：交互总量的主要效用损失来自结构性截断偏差 $E_{\text{clip}}$，而高斯噪声误差 $E_{\text{noise}}$ 在统计总量面前处于微弱数量级。
\end{techbox}

\subsection{任务成功率机制 (Success Rate Profile)}
\paragraph{1. 联合向量空间建模与共享截断：}
设用户 $u$ 的任务成功计数向量为 $s_u \in \mathbb{R}^D$，总调用计数向量为 $n_u \in \mathbb{R}^D$。显然有 $0 \le s_{u, j} \le n_{u, j}$。为了消除分子分母独立截断引发的逻辑失真（例如出现加噪后成功数反超总调用数的荒谬现象），系统设计了\textbf{共享比例截断算子}：
\begin{equation}
  a_u = \min\left(1, \frac{C_{\text{rate}}}{\|n_u\|_2}\right), \qquad \bar{n}_u = a_u n_u, \qquad \bar{s}_u = a_u s_u
\end{equation}
由于 $a_u \in (0, 1]$ 是标量，截断严格保留了用户个体的原始成功率比例关系（$\bar{s}_{u, j} / \bar{n}_{u, j} = s_{u, j} / n_{u, j}$）。将截断后的成功向量与总量向量拼接为 $2D$ 维联合向量：
\begin{equation}
  w_u = \begin{bmatrix} \bar{s}_u \\ \bar{n}_u \end{bmatrix} \in \mathbb{R}^{2D}
\end{equation}

\paragraph{2. 联合敏感度推导定理：}
\begin{techbox}{Success Rate 联合敏感度定理}
对于任意用户 $u$，联合向量 $w_u = [\bar{s}_u ; \bar{n}_u]$ 的 $L_2$ 范数满足：
\begin{equation}
  \|w_u\|_2 = \sqrt{\|\bar{s}_u\|_2^2 + \|\bar{n}_u\|_2^2} \le \sqrt{\|\bar{n}_u\|_2^2 + \|\bar{n}_u\|_2^2} = \sqrt{2}\|\bar{n}_u\|_2 \le \sqrt{2} C_{\text{rate}}
\end{equation}
因此，在用户级邻接关系下，$2D$ 维联合查询 $\sum_{u \in \mathcal{U}} w_u$ 的全局 $L_2$ 敏感度满足：
\begin{equation}
  \Delta_2 \le \sqrt{2} C_{\text{rate}}
\end{equation}
\end{techbox}

\paragraph{3. 单次联合释放与阈值抑制后处理：}
系统执行\textbf{单次 $2D$ 维高斯机制联合释放}：
\begin{equation}
  \tilde{W} = \begin{bmatrix} \tilde{S} \\ \tilde{N} \end{bmatrix} = \sum_{u \in \mathcal{U}} w_u + \mathcal{N}\left(0, \sigma^2 \mathbf{I}_{2D}\right)
\end{equation}
该操作仅产生一次隐私预算消耗。后处理算子严格定义如下：
\begin{enumerate}[label=\textbf{步骤\arabic*：}, leftmargin=4.5em]
  \item \textbf{总量非负约束：} $\hat{N}_j = \max(0, \tilde{N}_j)$；
  \item \textbf{分子物理约束：} $\hat{S}_j = \operatorname{clip}(\tilde{S}_j, 0, \hat{N}_j)$，确保成功数绝不超过总调用数；
  \item \textbf{小样本抑制判定：} 若 $\hat{N}_j < \tau$（默认阈值 $\tau = 5.0$），判定该单元加噪后样本量不足，可能引发除法剧烈扰动，系统返回 \texttt{SUPPRESSED / null}；
  \item \textbf{常规比率求取：} 若 $\hat{N}_j \ge \tau$，输出估计成功率 $\hat{r}_j = \hat{S}_j / \hat{N}_j \in [0.0, 1.0]$。
\end{enumerate}

\subsection{热门任务发现机制 (Top-K Tasks Profile)}
\paragraph{1. 有穷受管字典设定：}
明确系统的统计设定为\textbf{已知有限分类字典下的频率排序（Configured Finite Taxonomy）}，而非开放词表的重头项挖掘（Open-Universe Heavy Hitters）。设候选任务集合为 $\mathcal{T} = \{t_1, \dots, t_T\}$。

\paragraph{2. 用户参与度截断与敏感度：}
定义用户对任务的二值参与指示向量 $v_{u, t} = \mathbf{1}[m_{u, t} > 0]$。施加截断约束 $\|v_u\|_0 \le k_{\text{task}}$。聚合直方图在受管字典维度 $T$ 上的全局 $L_2$ 敏感度为：
\begin{equation}
  \Delta_2 = \sqrt{\min(k_{\text{task}}, T)} \quad (\text{代码保守校准使用 } \sqrt{k_{\text{task}}})
\end{equation}

\paragraph{3. 噪声注入与全字典排序合规准则：}
为了避免选择性数据窥探偏置，系统严格禁止“先按真实数据挑选前 $K$ 个任务再单独加噪”的不合规操作。合法流程必须先对字典中的全部 $T$ 个任务独立注入高斯噪声：
\begin{equation}
  \tilde{c}_t = \max\left(0, \sum_{u \in \mathcal{U}} v_{u, t} + \mathcal{N}(0, \sigma^2)\right), \quad \forall t \in \mathcal{T}
\end{equation}
随后仅基于加噪得分向量 $\tilde{c}$ 进行降序排序，取前 $K$ 个元素作为输出集合 $\operatorname{TopK}(\tilde{c})$。

\section{高斯差分隐私机制与参数精确校准}
\label{sec:gaussian}

\subsection{高斯机制的数学定义与特性}
高斯机制（Gaussian Mechanism）是差分隐私理论中处理实数向量查询的核心原语~\cite{dwork2006calibrating, dwork2014algorithmic}。设 $f: \mathcal{D} \to \mathbb{R}^d$ 为一个确定性向量查询函数，其全局 $L_2$ 敏感度定义为：
\begin{equation}
  \Delta_2(f) = \max_{D \sim D'} \|f(D) - f(D')\|_2
\end{equation}
高斯机制通过向真实向量 $f(D)$ 叠加各向同性多元高斯随机变量实现扰动：
\begin{equation}
  \mathcal{M}_{\text{Gauss}}(D) = f(D) + \mathcal{N}\left(0, \sigma^2 \mathbf{I}_d\right)
\end{equation}
高斯机制相较于拉普拉斯机制的关键理论优势在于：其噪声矢量范数增长与维度 $d$ 的解耦特性更契合高维统计需求，且在多重组合与长周期发布场景下，具有优良的 Rényi 散度可加性。

\subsection{经典高斯校准及其理论局限}
在早期差分隐私文献中，为了保证算法满足 $(\varepsilon, \delta)$-DP，通常采用基于切诺夫界的经典闭式近似公式：
\begin{equation}
  \label{eq:classical_gaussian}
  \sigma_{\text{classical}} = \frac{\Delta_2(f) \sqrt{2 \ln\left(\frac{1.25}{\delta}\right)}}{\varepsilon}
\end{equation}
然而，式~(\ref{eq:classical_gaussian}) 存在严格的理论适用边界与工程劣势：
\begin{itemize}[leftmargin=2em]
  \item \textbf{理论适用范围受限：} 该公式仅在 $\varepsilon \in (0, 1]$ 的窄区间内被严格证明满足定义；当 $\varepsilon > 1$ 时，该闭式界限不再构成严格的充分条件。
  \item \textbf{噪声方差过度放大（Over-conservatism）：} 经典充分条件公式基于高斯尾部放缩，计算得到的标准差 $\sigma$ 存在松弛。在样本量敏感的企业分析场景下，这种不必要的噪声膨胀将直接损害下游业务的决策可用性。
\end{itemize}

\subsection{解析高斯机制校准 (Analytic Gaussian Calibration)}
为了在保证数学证明严格性的同时最大化统计效用，本项目在底层机制中全面采用 Balle 与 Wang (ICML 2018) 提出的\textbf{解析高斯机制（Analytic Gaussian Mechanism）}~\cite{balle2018improving}。

\begin{techbox}{解析高斯校准定理 (Balle \& Wang, 2018)}
对于给定全局 $L_2$ 敏感度 $\Delta_2$ 以及目标差分隐私参数 $(\varepsilon, \delta)$，高斯扰动机制满足 $(\varepsilon, \delta)$-DP 当且仅当隐私损失随机变量的尾部累积概率函数满足：
\begin{equation}
  \label{eq:analytic_tail}
  \delta(\sigma; \varepsilon, \Delta_2) = \Phi\left(\frac{\Delta_2}{2\sigma} - \frac{\varepsilon\sigma}{\Delta_2}\right) - e^\varepsilon \Phi\left(-\frac{\Delta_2}{2\sigma} - \frac{\varepsilon\sigma}{\Delta_2}\right) \le \delta
\end{equation}
其中 $\Phi(x) = \frac{1}{\sqrt{2\pi}} \int_{-\infty}^x e^{-t^2/2} dt$ 为标准一维正态分布的累积分布函数（CDF）。
\end{techbox}

在工程实现中，系统将式~(\ref{eq:analytic_tail}) 的等号条件构造为关于噪声标准差 $\sigma$ 的一维严格单调函数：
\begin{equation}
  g(\sigma) = \Phi\left(\frac{\Delta_2}{2\sigma} - \frac{\varepsilon\sigma}{\Delta_2}\right) - e^\varepsilon \Phi\left(-\frac{\Delta_2}{2\sigma} - \frac{\varepsilon\sigma}{\Delta_2}\right) - \delta = 0
\end{equation}
采用 Brent 单变量求根算法（\texttt{scipy.optimize.}\allowbreak\texttt{brentq}），在初始包围区间 $[10^{-5}\Delta_2, \, \sigma_{\text{classical}} \times 2.5]$ 内求解精确根 $\sigma^*$，收敛绝对容差设为 $10^{-12}$。该机制不仅对任意 $\varepsilon > 0$ 与 $\delta \in (0, 1)$ 均全局成立，而且从数学上保证了在满足给定 $(\varepsilon, \delta)$ 保证下的高斯噪声方差达到理论极限最小值。

\subsection{实际参数校准与对比分析}
在基准安全配置（$\varepsilon = 1.0, \, \delta = 10^{-6}$）下，各统计指标的理论敏感度以及分别使用经典公式与解析高斯算法求得的噪声标度 $\sigma$ 对比如表~\ref{tab:sigma_comparison} 所示。

\begin{table}[htbp]
\centering
\small
\caption{基准参数配置下各统计指标的高斯噪声标度校准对比表 ($\varepsilon=1.0, \delta=10^{-6}$)}
\label{tab:sigma_comparison}
\begin{tabularx}{\textwidth}{lccccX}
\toprule
\textbf{统计指标} & \textbf{配置参数} & \textbf{$L_2$ 敏感度 $\Delta_2$} & \textbf{经典公式 $\sigma$} & \textbf{解析校准 $\sigma^*$} & \textbf{方差降低优化幅度} \\
\midrule
\textbf{User Prevalence} & $k = 9$ & $3.0000$ & $15.7513$ & $\mathbf{12.6740}$ & 噪声标准差降低 $19.53\%$ \\
\textbf{Normalized Share} & 天然有界 & $1.0000$ & $5.2504$ & $\mathbf{4.2247}$ & 噪声标准差降低 $19.54\%$ \\
\textbf{Interaction Volume} & $C = 20$ & $20.0000$ & $105.0084$ & $\mathbf{84.4936}$ & 噪声标准差降低 $19.54\%$ \\
\textbf{Success Rate} & $C_{\text{rate}} = 20$ & $28.2843$ & $148.5043$ & $\mathbf{119.4920}$ & 噪声标准差降低 $19.54\%$ \\
\textbf{Top-K Tasks} & $k_t = 5$ & $2.2361$ & $11.7403$ & $\mathbf{9.4467}$ & 噪声标准差降低 $19.54\%$ \\
\bottomrule
\end{tabularx}
\end{table}

在本文的基准安全配置（$\varepsilon = 1.0, \, \delta = 10^{-6}$）下，实测表明解析 Gaussian calibration 相对当前采用的经典充分条件公式将各指标的噪声标准差 $\sigma$ 降低约 $19.5\%$（噪声方差降低逾 $35.2\%$），显著增强了小样本聚合单元的抗畸变能力。在其他 $(\varepsilon, \delta)$ 参数设置下，降噪比例会随目标参数组合动态变动，系统不对其他参数组合做统一的百分比承诺。

\section{全生命周期隐私记账与发布控制机制}
\label{sec:accounting_release}

\subsection{Rényi 差分隐私 (RDP) 理论体系与紧致组合}
在企业实际运维中，统计发布是一项跨越数周乃至数月的长周期操作。根据差分隐私的基本组合定理，多次释放的 $\varepsilon$ 呈线性简单累加，会导致隐私预算迅速耗尽。为此，系统采用 Mironov (2017) 提出的 \textbf{Rényi 差分隐私（Rényi Differential Privacy, RDP）}~\cite{mironov2017renyi} 作为底层核心记账引擎。

\begin{definitionbox}{Rényi 差分隐私 (RDP)}
对于阶数 $\alpha > 1$，设两个概率分布 $P, Q$ 在同一支集上的 Rényi 散度定义为 $D_\alpha(P \parallel Q) = \frac{1}{\alpha - 1} \ln \int P(x)^\alpha Q(x)^{1-\alpha} dx$。若随机化机制 $\mathcal{M}$ 满足对于任意邻近数据集 $D \sim D'$：
\begin{equation}
  D_\alpha\left(\mathcal{M}(D) \,\parallel\, \mathcal{M}(D')\right) \le \rho(\alpha)
\end{equation}
则称机制 $\mathcal{M}$ 满足 $(\alpha, \rho(\alpha))$-RDP。
\end{definitionbox}

高斯机制在 RDP 框架下具备解析封闭的散度表达式。对于标准差为 $\sigma$、全局 $L_2$ 敏感度为 $\Delta_2$ 的单次高斯释放，其在任意阶数 $\alpha$ 上的 RDP 消耗量为：
\begin{equation}
  \label{eq:rdp_cost}
  \rho(\alpha) = \frac{\alpha \Delta_2^2}{2 \sigma^2}
\end{equation}
RDP 最大的理论与工程优势在于其在独立组合下的\textbf{严格线性可加性}：设系统先后执行了 $m$ 次具有不同敏感度与噪声标度的高斯发布，则系统在阶数 $\alpha$ 上的全生命周期累积隐私损失为单次损失的直接求和：
\begin{equation}
  \rho_{\text{life}}(\alpha) = \sum_{i=1}^m \rho_i(\alpha) = \sum_{i=1}^m \frac{\alpha \Delta_{2, i}^2}{2 \sigma_i^2}
\end{equation}
当需要向业务层或审计合规方呈现等价的 $(\varepsilon, \delta)$ 状态时，系统遍历标准离散阶数集合 $\mathcal{A} = [1.25, 1.5, 2.0, 3.0, 4.0, 5.0, 8.0, 10.0, 16.0, 32.0, 64.0, 128.0]$，执行精确的对偶凸优化转换：
\begin{equation}
  \label{eq:rdp_conversion}
  \varepsilon_{\text{life}}(\delta) = \min_{\alpha \in \mathcal{A}} \left[ \rho_{\text{life}}(\alpha) + \frac{\ln(1/\delta)}{\alpha - 1} \right]
\end{equation}
相较于经典高级组合定理（Advanced Composition Theorem），RDP 组合消除了复杂的松弛常数项，显著延长了企业在给定隐私上限下的有效可用发布次数。

\subsection{全生命周期隐私预算硬控制 (Lifetime Budget Cap)}
企业必须防范内部发布频次失控导致的隐私保证实质性失效。系统确立了“全生命周期终身隐私预算”（Lifetime Privacy Budget）控制规范：
\begin{enumerate}[label=(\alph*), leftmargin=2.5em]
  \item \textbf{政策配置上限 ($\varepsilon_{\text{cap}}$)：} 管理层或数据安全团队根据法规要求配置终身预算硬上限（系统默认 $\varepsilon_{\text{cap}} = 4.0, \, \delta_{\text{global}} = 10^{-6}$）。
  \item \textbf{状态持久化原则（Persistent State）：} 历史发布的 RDP 账本存储于持久化、事务一致的生命周期隐私账本中；正常服务重启、系统升级或前端刷新不会自动重置历史状态。
  \item \textbf{原子比较与预留（Atomic Compare-and-Reserve）：} 在每次发起实际计算与加噪之前，记账引擎首先模拟计算若本次发布提交后，转换得到的全生命周期 $\varepsilon_{\text{projected}}(\delta_{\text{global}})$。若满足：
    \begin{equation}
      \varepsilon_{\text{projected}}(\delta_{\text{global}}) > \varepsilon_{\text{cap}}
    \end{equation}
    系统立即抛出 \texttt{PrivacyBudgetExceededError}，无条件熔断并拒绝执行，坚决恪守 Fail-Closed 安全底线。
\end{enumerate}

\subsection{接收方作用域敏感组合 (Recipient-Aware Composition)}
针对多模型供应商协同反馈场景，若将所有外部接收方与内部审计视同单一主体，会导致内部高频评测迅速消耗外部共享预算；反之，若为每个外部供应商分配完全独立的无限预算，则在供应商合谋时将引发隐私防线的被击穿。系统研发了\textbf{接收方作用域敏感记账体系}，支持三种核心管理范式：
\begin{itemize}[leftmargin=2em]
  \item \texttt{GLOBAL\_}\allowbreak\texttt{EXTERNAL}（默认外部模式）：所有不可信外部供应商（OpenAI、Kimi、DeepSeek）共享同一账本作用域 $\mathcal{S}_{\text{ext}}$。任何向供应商 A 的发布，其隐私成本均在供应商 B 的可用预算中扣减，从而形式化抵御了多供应商联合反演攻击。
  \item \texttt{PER\_RECIPIENT}：在确立了严格商业互信且拥有相互审计与防泄密法律契约的特定实体之间，启用独立账本记账。
  \item \texttt{COLLUSION\_GROUP}：显式定义指定合谋组（Collusion Group），组内实体预算互通，组间相互隔离。
\end{itemize}

\subsection{规范发布与幂等重放机制 (Canonical Release \& Replay)}
\paragraph{均值恢复攻击风险：}
若外部接收方对完全相同的统计问题（例如连续多次查询“上月 OpenAI 编程任务成功率”）重复发起请求，且系统在每次请求时均独立采样高斯噪声返回，攻击者只需发起 $K$ 次查询并取均值：
\begin{equation}
  \bar{Y} = \frac{1}{K} \sum_{k=1}^K \left( f(D) + Z_k \right) = f(D) + \frac{1}{K} \sum_{k=1}^K Z_k \sim \mathcal{N}\left(f(D), \, \frac{\sigma^2}{K} \mathbf{I}\right)
\end{equation}
当 $K \to \infty$ 时，噪声方差以 $O(1/K)$ 速率衰减，差分隐私提供的稳定性保障显著削弱。

\paragraph{规范语义哈希与零成本重放：}
为了有效防范该特定攻击路径，系统引入了\textbf{规范发布引擎（Canonical Release Engine）}。系统对查询请求的所有本征要素提取确定性哈希签名：
\begin{equation}
  \text{ReleaseKey} = \text{SHA256}\Big(\text{fp}_{\text{dataset}} \,\|\, \text{window} \,\|\, \text{metric\_id} \,\|\, \mathcal{S}_{\text{scope}} \,\|\, \varepsilon \,\|\, \delta \,\|\, \text{taxonomy\_ver}\Big)
\end{equation}
查询执行时遵循如图~\ref{fig:flowchart} 所示的状态检查逻辑：
\begin{itemize}[leftmargin=2em]
  \item 若 \texttt{ReleaseKey} 在持久化发布库中已存在对应已提交（COMMITTED）记录，系统直接提取已有发布结果返回，标记 \texttt{cache\_hit = True}，且本次请求的\textbf{净增隐私成本恒为零 ($\Delta\text{RDP} = 0$)}；
  \item 仅当 \texttt{ReleaseKey} 为全系统首次出现时，才由操作系统提供的熵源（OS-provided entropy）初始化的随机数生成器生成单次高斯样本并持久化归档。
\end{itemize}
该机制从工程体系上确保了相同语义的分析任务在任意高频调用下，均无法被外部实体用于统计均值恢复，同时保障了前端控制台反复刷新时的体验与零额外预算支出。

\section{流式时间序列持续观测机制 (Continual Observation)}
\label{sec:continual}

\subsection{长周期动态监控的隐私与效用挑战}
在企业大模型日常治理中，管理团队不仅需要离线的单次汇总报表，更需要连续跟踪使用趋势。例如监控过去 7 天或 14 天的滑动窗口交互量波动，以便及时捕捉因业务系统变更引发的调用激增或异常下跌。

若将每一天的增量调用简单视为一次独立的差分隐私释放任务，连续发布 $T$ 天将导致两个无法克服的缺陷：
\begin{enumerate}[label=(\arabic*), leftmargin=2em]
  \item \textbf{隐私预算迅速耗尽：} 依据组合定理，独立加噪的 RDP 散度与发布天数 $T$ 呈严格线性增长（$\rho_{\text{total}} = O(T)$）。在一个季度（$T=90$）内，全生命周期预算上限将被迅速穿透。
  \item \textbf{滑动窗口方差累加：} 评估一个长度为 $W$ 天的滑动窗口总量时，若由 $W$ 个独立加噪的单日数据求和得到，噪声方差将线性叠加至 $W \sigma^2$，导致区间趋势严重失真。
\end{enumerate}

针对此问题，系统将 Dwork 等人 (STOC 2010)~\cite{dwork2010continual} 与 Chan 等人 (SODA 2011)~\cite{chan2011private} 开创的\textbf{差分隐私时间二叉树机制（Binary Time Tree）}工程化引入本项目，重点赋能交互总量（Interaction Volume）的持续观测。

\subsection{双层贡献截断体系 (Two-Tier Bounding)}
在持续时间序列中，单个员工不仅在单日可能发起极端调用，还可能在整月内持续高频消耗。为了建立有限全局敏感度，系统构建了\textbf{双层贡献截断体系}：
\begin{itemize}[leftmargin=2em]
  \item \textbf{第一层：单日贡献截断 ($C_{\text{day}}$)：}
    对用户 $u$ 在日期 $t$ 的单日交互计数向量 $y_{u, t} \in \mathbb{R}^D$ 实施局部 $L_2$ 范数截断：
    \begin{equation}
      \bar{y}_{u, t} = y_{u, t} \cdot \min\left(1, \frac{C_{\text{day}}}{\|y_{u, t}\|_2}\right)
    \end{equation}
  \item \textbf{第二层：代际综合贡献截断 ($C_{\text{epoch}}$)：}
    在预设的代际周期（Epoch Horizon，如 $T_{\text{epoch}} = 16$ 或 $30$ 天）内，约束员工在该周期内全部日期的截断向量合成范数：
    \begin{equation}
      \sqrt{\sum_{t=1}^{T_{\text{epoch}}} \|\bar{y}_{u, t}\|_2^2} \le C_{\text{epoch}}
    \end{equation}
    若超标，则执行全局等比收缩，确保该员工在整个代际时间树上的综合 $L_2$ 敏感度严格受限于 $C_{\text{epoch}}$。
\end{itemize}

\subsection{时间二叉树架构与规范二进节点划分}
系统将一个代际内的连续时间区间映射为一棵规范时间二叉树（Binary Time Tree）。以长度为 8 天的周期为例，其节点层级拓扑如图~\ref{fig:binary_tree} 所示。

\begin{figure}[htbp]
\centering
\begin{tikzpicture}[
  scale=0.88, every node/.style={scale=0.88},
  treenode/.style={draw=navyblue, fill=lightbg, rounded corners=3pt, thick, align=center, inner sep=4pt, font=\small},
  leafnode/.style={draw=slateblue, fill=white, rounded corners=2pt, thick, align=center, inner sep=3pt, font=\footnotesize},
  querynode/.style={draw=accentgreen, fill=accentgreen!15, rounded corners=3pt, very thick, align=center, inner sep=4pt, font=\small},
  edge from parent/.style={draw=navyblue, thick, -{Latex[length=1.5mm]}},
  level distance=1.3cm,
  sibling distance=3.6cm
]

  % Root (Level 3)
  \node[treenode] (root) at (0, 0) {节点 [1, 8]\\(Level 3)}
    child {node[querynode] (n14) {节点 [1, 4]\\(Level 2)}
      child {node[treenode] (n12) {[1, 2]\\(L1)}
        child {node[leafnode] (d1) {第1天\\(L0)}}
        child {node[leafnode] (d2) {第2天\\(L0)}}
      }
      child {node[treenode] (n34) {[3, 4]\\(L1)}
        child {node[leafnode] (d3) {第3天\\(L0)}}
        child {node[leafnode] (d4) {第4天\\(L0)}}
      }
    }
    child {node[treenode] (n58) {节点 [5, 8]\\(Level 2)}
      child {node[querynode] (n56) {节点 [5, 6]\\(Level 1)}
        child {node[leafnode] (d5) {第5天\\(L0)}}
        child {node[leafnode] (d6) {第6天\\(L0)}}
      }
      child {node[treenode] (n78) {[7, 8]\\(L1)}
        child {node[querynode] (d7) {第7天\\(Level 0)}}
        child {node[leafnode] (d8) {第8天\\(L0)}}
      }
    };

\end{tikzpicture}
\caption{时间二叉树（Binary Time Tree）结构与区间规范分解示意图 (查询区间 [1, 7] 仅由 3 个加噪节点精确合成：$\text{Node}_{[1,4]} + \text{Node}_{[5,6]} + \text{Node}_{[7,7]}$)}
\label{fig:binary_tree}
\end{figure}

时间二叉树具备如下核心数学特性：
\begin{enumerate}[label=(\alph*), leftmargin=2.5em]
  \item \textbf{层级正交性（Level Orthogonality）：} 树的每一层（Level）由对整个时间跨度的不重叠二进区间划分构成。因此，在任意单一层级内部，每个员工的数据至多仅落入该层级的一个节点中。
  \item \textbf{对数级区间分解（Logarithmic Decomposition）：} 对于任意合法的历史时间区间 $[t_{\text{start}}, t_{\text{end}}]$（长度为 $W$），系统利用标准的二进区间分解算法，将其唯一分解为\textbf{至多 $2 \lfloor \log_2 W \rfloor$ 个互不相交的预置二叉树节点}。
\end{enumerate}

\subsection{持久化加噪节点与抗均值攻击}
在工程实现中，时间二叉树中的每个二进节点在首度被物化时，注入标准差为 $\sigma_{\text{tree}}$ 的独立高斯噪声，随后持久化写入节点数据库中。
系统确立了严格的作用域节点复用准则：\textbf{在同一隐私组合作用域（Privacy Composition Scope）内，对相同规范节点（Canonical Node）复用相同持久化加噪值（Persistent Noisy Value）}。具体行为遵循以下作用域隔离矩阵：
\begin{itemize}[leftmargin=2em]
  \item \textbf{相同 \texttt{GLOBAL\_EXTERNAL} 作用域：} 处于该共享作用域下的外部供应商（如 OpenAI、Kimi）请求重叠区间时，严格复用已物化的持久化加噪节点，在内存中进行线性加和，避免生成独立重复噪声样本，防范特定重复查询平均攻击；
  \item \textbf{相同 \texttt{COLLUSION\_GROUP} 作用域：} 显式划分在同一合谋组内的接收方，共享并复用相同节点加噪值与关联预算池；
  \item \textbf{不同 \texttt{COLLUSION\_GROUP} 与独立 \texttt{PER\_RECIPIENT} 作用域：} 依据系统配置与隔离实现进行独立验证与记账，在不同作用域间实施严格的上下文隔离，不得无条件跨独立作用域共享同一节点噪声。
\end{itemize}

\subsection{前缀稳定性与代际边界管理}
持续观测系统在生产环境必须遵守严格的因果性准则：
\begin{itemize}[leftmargin=2em]
  \item \textbf{前缀稳定性（Prefix Stability）：} 随着未来日期的推移，已完成闭环封存的过去历史节点的数值不得发生任何变动。
  \item \textbf{历史归档边界（Finalized Frontier）：} 系统引入封存时间边界 $T_{\text{finalized}}$。凡落后于该边界的迟到数据（Late-arriving Records），系统依据 Fail-Closed 策略严禁直接写回已有已发布的历史节点，防范历史区间的隐私保证界被隐式击穿。
  \item \textbf{跨代际预算累积（Epoch Boundary）：} 当一个代际（如第 30 天）结束并开启下一个代际树时，系统明确\textbf{代际翻转绝不意味着隐私预算归零清空}。新代际二叉树的 RDP 散度无缝累加至全系统统一的持久化账本中，确保终身预算受控。
\end{itemize}

\section{差分隐私效用与精度实证评测}
\label{sec:utility}

\subsection{效用度量定义与两层误差分解规范}
在评估差分隐私发布算法的工程可用性时，必须建立严谨的误差基准。本系统对统计过程中的三个核心状态建立明确的数学定义：
\begin{enumerate}[label=(\alph*), leftmargin=2.5em]
  \item \textbf{原始精确值 ($M_{\text{raw}}$)：} 在完全未施加任何隐私干预（无截断、无噪声）的前提下，基于内部真实遥测明细计算得到的数学真值；
  \item \textbf{截断精确值 ($M_{\text{clip}}$)：} 经过用户级贡献限制算子（如 $L_0$ 单元截断、单纯形映射或 $L_2$ 范数截断）处理后、但在注入差分隐私高斯噪声之前计算得到的中间真值；
  \item \textbf{差分隐私估计值 ($M_{\text{DP}}$)：} 叠加高斯噪声并执行合法后处理（如值域截断、归一化与小样本抑制）后，最终面向接收方发布的脱敏估计值。
\end{enumerate}

系统确立了\textbf{两层误差解耦分析准则}：
\begin{equation}
  \label{eq:two_tier_decomp}
  E_{\text{total}} = M_{\text{DP}} - M_{\text{raw}} = \underbrace{(M_{\text{clip}} - M_{\text{raw}})}_{E_{\text{clip}}\text{ (结构性截断偏差)}} + \underbrace{(M_{\text{DP}} - M_{\text{clip}})}_{E_{\text{noise}}\text{ (差分隐私加噪误差)}}
\end{equation}
针对不同业务指标，系统采用针对性的客观误差度量指标：
\begin{itemize}[leftmargin=2em]
  \item \textbf{百分点绝对误差（Percentage-Point Error, pp）：} 适用于 Prevalence 与 Success Rate 等比率指标，定义为 $E_{\text{pp}} = |\hat{r} - r_{\text{raw}}| \times 100$；
  \item \textbf{平均相对误差（Mean Relative Error, MRE）：} 适用于 Interaction Volume 等总量计数指标，定义为 $\text{MRE} = \frac{|\hat{M} - M_{\text{raw}}|}{\max(|M_{\text{raw}}|, 1.0)}$；
  \item \textbf{曼哈顿距离（$L_1$ Distance）：} 适用于 Normalized Share 概率分布向量，定义为 $L_1(s_{\text{raw}}, \hat{s}) = \frac{1}{2} \sum_{j=1}^D |s_{\text{raw}, j} - \hat{s}_j|$；
  \item \textbf{Top-$K$ 命中率与 Jaccard 指数：} 适用于排序挖掘指标，评估候选集合与真实排名的重合比例。
\end{itemize}

\subsection{用户使用率 (Prevalence) 实测与 92.85\% 根因剖析}
在基准企业环境（在册员工 $N_{\text{public}} = 50,000$ 人、包含 3,011,232 条真实交互记录的数据集 \texttt{interactions\_n50000.parquet}，安全配置 $\varepsilon = 1.0, \delta = 10^{-6}$）下，系统对全部 9 个细分业务单元执行了全量实测，结果如表~\ref{tab:prevalence_benchmark} 所示。

\begin{table}[htbp]
\centering
\small
\caption{用户使用率 (User Prevalence) 50,000 人基准实测表 ($\varepsilon=1.0, \delta=10^{-6}, N=50,000$)}
\label{tab:prevalence_benchmark}
\resizebox{\textwidth}{!}{%
\begin{tabular}{llccccc}
\toprule
\textbf{模型供应商} & \textbf{业务任务} & \textbf{真实用户数} & \textbf{原始使用率} & \textbf{截断使用率} & \textbf{DP 估计值} & \textbf{百分点误差 (pp)} \\
\midrule
OpenAI (\texttt{provider\_a}) & Technical & $46,423$ & $92.85\%$ & $92.85\%$ & $92.85\%$ & $+0.01\text{ pp}$ \\
OpenAI (\texttt{provider\_a}) & Document & $35,941$ & $71.88\%$ & $71.88\%$ & $71.86\%$ & $-0.03\text{ pp}$ \\
OpenAI (\texttt{provider\_a}) & Analysis & $37,374$ & $74.75\%$ & $74.75\%$ & $74.77\%$ & $+0.02\text{ pp}$ \\
Kimi (\texttt{provider\_b}) & Technical & $38,451$ & $76.90\%$ & $76.90\%$ & $76.93\%$ & $+0.02\text{ pp}$ \\
Kimi (\texttt{provider\_b}) & Document & $46,379$ & $92.76\%$ & $92.76\%$ & $92.71\%$ & $-0.05\text{ pp}$ \\
Kimi (\texttt{provider\_b}) & Analysis & $37,543$ & $75.09\%$ & $75.09\%$ & $75.05\%$ & $-0.03\text{ pp}$ \\
DeepSeek (\texttt{provider\_c}) & Technical & $35,772$ & $71.54\%$ & $71.54\%$ & $71.55\%$ & $+0.00\text{ pp}$ \\
DeepSeek (\texttt{provider\_c}) & Document & $38,483$ & $76.97\%$ & $76.97\%$ & $76.96\%$ & $-0.01\text{ pp}$ \\
DeepSeek (\texttt{provider\_c}) & Analysis & $45,734$ & $91.47\%$ & $91.47\%$ & $91.47\%$ & $+0.00\text{ pp}$ \\
\midrule
\multicolumn{6}{l}{\textbf{全局平均百分点绝对误差 (Mean Absolute pp Error)：}} & $\mathbf{0.02\text{ pp}}$ \\
\bottomrule
\end{tabular}%
}
\end{table}

\paragraph{92.85\% 数值事实的深入调查与澄清：}
在前期内部评审中，曾有审查意见提出疑虑：“页面显示 OpenAI 编程辅助使用率仅约 92.85\%，若全员使用率接近 100\%，是否存在系统导致了约 7 个百分点的异常下降？”

经本期审计建立的端到端数值回溯，得出了确定性事实澄清：
\begin{enumerate}[label=(\arabic*), leftmargin=2em]
  \item \textbf{统计口径的混淆：} 在册的 50,000 名员工中，至少调用过大模型一次的员工总数确实为 50,000 人（全公司宏观使用率为 100.0\%）；曾使用过 OpenAI 处理任意任务的员工为 49,592 人（占 99.18\%）。但是，\textbf{具体到“使用 OpenAI 处理编程辅助（Technical）”这一特定单元的独立员工数客观上恰好为 46,423 人}。
  \item \textbf{数据事实结论：} 原始真实使用率即为 $46,423 / 50,000 = \mathbf{92.846\% \approx 92.85\%}$。7 个百分点的差距是企业员工客观工作分工的反映（有 3,577 名员工在观测周期内仅使用大模型撰写文档或分析报表，从不编写代码）。
  \item \textbf{差分隐私极微误差：} 高斯噪声在此处的实际扰动仅为 $+3.86$ 人，折合百分点误差仅为 $\mathbf{+0.01\text{ pp}}$。该实测有力推翻了“差分隐私导致数据剧烈失真”的常见误解。
\end{enumerate}

\subsection{任务成功率 (Success Rate) 实测结果}
任务成功率机制在共享截断算子（$C_{\text{rate}} = 20$）与联合释放设计下展现出极高的可用性。在 1,000 用户与 50,000 用户基准下的实测结果如表~\ref{tab:success_rate_benchmark} 所示。

\begin{table}[htbp]
\centering
\small
\caption{任务成功率 (Success Rate) 实测精度对比表 ($\varepsilon=1.0, \delta=10^{-6}$)}
\label{tab:success_rate_benchmark}
\resizebox{\textwidth}{!}{%
\begin{tabular}{lcccccc}
\toprule
\textbf{测试数据集} & \textbf{真实成功数} & \textbf{真实调用量} & \textbf{原始成功率} & \textbf{截断成功率} & \textbf{DP 估计值} & \textbf{百分点误差 (pp)} \\
\midrule
$N=1,000$ (合成基准) & $5,434$ & $6,396$ & $84.96\%$ & $84.88\%$ & $86.52\%$ & $+1.56\text{ pp}$ \\
$N=50,000$ (企业基准) & $272,118$ & $320,182$ & $84.99\%$ & $84.99\%$ & $85.08\%$ & $\mathbf{+0.09\text{ pp}}$ \\
\bottomrule
\end{tabular}%
}
\end{table}

实测表明，在万级员工规模下，高斯机制注入的噪声相较于数十万级的真实样本量影响微弱，成功率估计误差迅速收敛至 $0.09\text{ pp}$（相对误差小于 $0.11\%$），完全满足采购评估与 SLA 质量仲裁的商业需求。

\subsection{用户归一化占比 (Normalized Share) 实测结果}
在 50,000 用户基准下，用户工作流结构分布的实测对比表明：各业务单元的真实占比位于 $7.25\%$ 至 $20.66\%$ 之间，DP 估计值与真实值的平均百分点误差为 $\mathbf{0.015\text{ pp}}$，单纯形曼哈顿距离 $L_1 < 0.0003$，输出严格满足和为 1 的概率分布性质。

\subsection{交互总量 (Interaction Volume) 评测与截断偏差归因}
交互总量统计的实测数据如表~\ref{tab:volume_benchmark} 所示，该结果深刻揭示了长尾数据在差分隐私截断下的本质挑战。

\begin{table}[htbp]
\centering
\small
\caption{交互总量 (Interaction Volume) 两层误差解耦实测数据表 ($C=20, \varepsilon=1.0, \delta=10^{-6}$)}
\label{tab:volume_benchmark}
\resizebox{\textwidth}{!}{%
\begin{tabular}{lcccccc}
\toprule
\textbf{数据规模} & \textbf{原始真实值 $M_{\text{raw}}$} & \textbf{截断真值 $M_{\text{clip}}$} & \textbf{DP 发布值 $M_{\text{DP}}$} & \textbf{截断偏差 $E_{\text{clip}}$} & \textbf{加噪误差 $E_{\text{noise}}$} & \textbf{综合相对误差} \\
\midrule
$N=1,000$ & $12,380$ & $7,606.6$ & $7,632.4$ & $-4,773.4$ ($-38.56\%$) & $+25.7$ ($+0.21\%$) & $38.35\%$ \\
$N=50,000$ & $622,241$ & $374,208$ & $374,120$ & $-248,033$ ($-39.86\%$) & $-88$ ($-0.02\%$) & $\mathbf{39.88\%}$ \\
\bottomrule
\end{tabular}%
}
\end{table}

\paragraph{效用归因的客观技术结论：}
在两组测试中，综合相对误差均接近 $40\%$。但通过系统的两层误差解耦公式：
\begin{itemize}[leftmargin=2em]
  \item 差分隐私高斯加噪产生的随机误差 $E_{\text{noise}}$ 仅为 $-88$ 次，对总量的相对扰动仅为 $\mathbf{0.02\%}$；
  \item \textbf{超过 99.9\% 的效用损失来自于预设的常数贡献截断阈值 $C=20$。} 在真实业务中，部分自动化脚本或重度研发人员单月调用成百上千次，其超过 20 次的部分全部被几何截断。
\end{itemize}
这一发现明确了系统的架构边界：\textbf{交互总量的误差瓶颈不在于高斯机制，而在于静态截断参数与重度用户调用长尾之间的结构性矛盾}。

\subsection{热门任务 (Top-K Tasks) 挖掘效用}
在当前预置的 3 任务分类体系下，真实调用量排名前三位依次为：Technical (48,432 人)、Document (46,379 人)、Analysis (45,734 人)。系统在注入高斯噪声（$\sigma = 9.45$）后，排序位次完全保持一致，Top-3 命中率达到 $100\%$，Jaccard 相似度为 $1.00$。系统在报告中客观注明：当前受管字典基数相对有限，该指标的高命中率存在字典规模较小的客观前提。

\subsection{隐私预算 \texorpdfstring{$\varepsilon$}{Epsilon} 与统计效用敏感性扫描 (\texorpdfstring{$\varepsilon$}{Epsilon}-Sweep)}
为了给数据治理团队制定合理的发布政策提供实证支撑，系统在隔离实验环境执行了多轮参数扫描（$\varepsilon \in [0.1, 4.0]$，每点运行 10 次独立随机试验取均值）。实测趋势表明：
\begin{itemize}[leftmargin=2em]
  \item 比率类指标（Prevalence 与 Success Rate）在 $\varepsilon \ge 0.5$ 时即已收敛至 $0.2\text{ pp}$ 以内的极高精度区间，进一步放大 $\varepsilon$ 对可用性的边际改善递减；
  \item 当 $\varepsilon < 0.25$ 时，高斯噪声标度剧增，小样本单元的被抑制率显著上升。因此，系统将默认单次查询预算推荐设定为 $\varepsilon = 1.0$。
\end{itemize}

\section{工程实现与安全控制架构}
\label{sec:engineering}

\subsection{生产环境真随机数控制与注入防御}
差分隐私的数学安全性从根本上依赖于噪声采样的不可预测性。若外部攻击者能够预测、操纵或重放随机数生成器（RNG）的种子（Seed），差分隐私保证将完全坍塌。

针对该安全风险，系统在代码架构层实施了严格的\textbf{调用方 RNG 控制阻断机制}：
\begin{itemize}[leftmargin=2em]
  \item \textbf{生产模式熵源强制锁定：} 在生产发布接口 \texttt{AnalyticsService.compute()} 中，系统无条件采用基于操作系统提供的熵源（OS-provided entropy）初始化的随机数生成器（\texttt{np.random.default\_rng()}）；
  \item \textbf{参数探查与即时熔断：} 任何调用方若尝试传入包含 \texttt{rng}、\texttt{seed}、\texttt{noise\_seed}、\texttt{noise} 或 \texttt{random\_state} 的参数，代码拦截器立即触发 \texttt{ProductionRNGControlError} 并阻断执行；
  \item \textbf{禁止动态覆盖截断界：} 严禁外部调用方在查询时动态覆写截断阈值 \texttt{clip\_norm}，所有截断边界必须且仅能在受管的中心化配置中统一固化。
\end{itemize}

\subsection{评测实验与生产环境状态严格隔离}
为了支持算法团队开展长期的效用基准评测，同时坚决避免将测试流量混入具有法律合规效力的正式账本，系统构建了\textbf{双模态环境隔离执行体系}：
\begin{itemize}[leftmargin=2em]
  \item \textbf{生产服务管线（Production Pipeline）：} 仅响应标准发布请求，严格受控于全生命周期预算上限，所有已提交记录持久化写入生产数据库 \texttt{results/privacy\_state.db}。
  \item \textbf{实验评测管线（Experiment Runner）：} 通过专用内部通道 \texttt{compute\_experiment()} 调用，支持传入确定性种子以便实验复现，其内部记账被重定向至隔离的临时数据库 \texttt{results/experiment\_state.db}。系统部署了防护守卫（Fencing Guard）：若在实验模式下检测到数据库路径意外指向生产文件，系统将拒绝启动。
\end{itemize}

\subsection{关系型持久化状态机与并发事务控制}
为了保障在多业务部门并发查询场景下的记账一致性，系统依托轻量级关系型数据库与预写日志（WAL）机制，实现了\textbf{具备租约隔离与世代栅栏（Generation Fencing）的原子预留状态机}。状态转移逻辑包含以下阶段：
\begin{enumerate}[label=(\alph*), leftmargin=2.5em]
  \item \textbf{预留阶段（RESERVED）：} 开启数据库排他事务，核查历史累积 RDP 向量；若未超标，向账本插入处于租约有效期（\texttt{lease\_timeout = 300s}）内的预留令牌，并锁定候选预算；
  \item \textbf{计算与加噪阶段（EXECUTING）：} 在事务外异步执行矩阵聚合与高斯机制采样，避免长耗时计算阻塞全局数据库锁；
  \item \textbf{提交阶段（COMMITTED）：} 携带预留世代编号（Generation Token）执行二次原子事务提交；若工作节点在此期间发生崩溃或超时，未完成提交的预留令牌将被后续回收线程安全回收，防止预算悬空泄露或重复双花（Double-Spending）。
\end{enumerate}

\subsection{全面防御性故障闭环机制 (Fail-Closed Architecture)}
系统确立了“任何异常无条件阻断、绝不静默放行、绝不降级明文”的\textbf{全闭环安全准则（Fail-Closed）}：
\begin{itemize}[leftmargin=2em]
  \item \textbf{模式缺失拦截：} 成功率查询若遭遇缺失 \texttt{success} 标记的遥测日志，抛出 \texttt{MissingRequired}\allowbreak\texttt{StatisticColumn} 异常，拒绝执行；
  \item \textbf{分母失配拦截：} 上传数据中的活跃独立用户数若超过声明的公开在册人口（$n_{\text{active}} > N_{\text{public}}$），抛出 \texttt{DatasetConfiguration}\allowbreak\texttt{MismatchError}，阻断非法的分母越界；
  \item \textbf{越权窃取拦截：} 任何外部不可信接收方尝试请求精确对比报表（\texttt{include\_exact = True}）或访问内部输入预览接口，系统统一返回 HTTP 403 强安全违例（\texttt{RecipientSecurity}\allowbreak\texttt{Violation}）。
\end{itemize}

\subsection{多端发布数据一致性保证}
为了避免前后端在百分比转换、浮点精度舍入或接收方映射时产生不一致，系统确立了\textbf{“单一权威事实源（Single Source of Truth）”}架构：
\begin{center}
  \textbf{数据库已提交记账快照} $\equiv$ \textbf{RESTful API JSON 载荷} $\equiv$ \textbf{Web 控制台呈现视图} $\equiv$ \textbf{CSV / JSON 导出文件}
\end{center}
前端 Web 界面被严格限制为只读视图呈现器（View-only Presenter），禁止在浏览器端执行任何隐式的数学重加噪或二次缩放。CSV/JSON 文件的导出直接流式序列化数据库中已提交的快照记录，\textbf{下载操作不产生任何随机数重抽样，也不增加任何隐私预算消耗}。

\section{独立验证与工程测试基准}
\label{sec:validation}

\subsection{全量回归测试套件执行基准}
本项目构建了严密的自动化多层次测试套件，全面覆盖底层代数算子、机制校准、并发数据库事务、RDP 记账账本、外部安全隔离以及 Web 控制台 REST 接口。在当前开发分支（\texttt{v1.1-dev}）上，执行全量自动化回归测试的结果如表~\ref{tab:test_summary} 所示。

\begin{table}[htbp]
\centering
\small
\caption{自动化测试工程回归验证基准结果 (Python 3.14 环境)}
\label{tab:test_summary}
\begin{tabularx}{\textwidth}{lccX}
\toprule
\textbf{测试分类与子系统} & \textbf{测试用例数} & \textbf{执行状态} & \textbf{重点覆盖的安全与算法属性} \\
\midrule
\texttt{test\_statistic\_profiles.py} & $28$ & 全部 PASS & 5 大指标专属 Profile 数学规范与后处理单调性 \\
\texttt{test\_profile\_accounting.py} & $14$ & 全部 PASS & 单次释放 RDP 描述子与联合发布敏感度合法性 \\
\texttt{test\_continual\_*.py} (7个模块) & $82$ & 全部 PASS & 时间二叉树节点分解、跨接收方节点复用与最终边界 \\
\texttt{test\_cache\_and\_isolation.py} & $12$ & 全部 PASS & 规范重放幂等性、外部禁止精确值访问与供应商隔离 \\
\texttt{test\_privacy\_state.py} & $36$ & 全部 PASS & SQLite 原子事务、世代栅栏、租约超时与并发争用 \\
\texttt{test\_web\_*.py} (11个模块) & $94$ & 全部 PASS & Web 控制台 API、KaTeX 渲染、CSV 导出与策略更新 \\
其余机制、规划器与集成测试 & $276$ & 全部 PASS & 摄取校验、高斯校准、动态分类法与长周期累加 \\
\midrule
\textbf{全量汇总 (Total Baseline)} & $\mathbf{542}$ & $\mathbf{0\text{ Failed}}$ & $\mathbf{542\text{ 项测试全量通过}}$ (总耗时 $573.10\text{ 秒}$) \\
\bottomrule
\end{tabularx}
\end{table}

\subsection{敏感度相邻数据集 (\texorpdfstring{$D \sim D'$}{D ~ D'}) 独立构造验证}
为了防止代码中声称的敏感度与算法实际聚合产生的位移脱节，诊断套件独立构建了边界相邻数据集验证程序。测试构造了包含空用户、单单元极值调用（1,000 次）、全单元均匀调用以及全单元极端暴增（10,000 次）在内的 6 种极端用户向量 $u^*$，独立计算 $\|f(D) - f(D')\|_2$。实测最大经验位移与代码声明敏感度对比如表~\ref{tab:sensitivity_verification} 所示。

\begin{table}[htbp]
\centering
\small
\caption{五大统计指标相邻数据集经验敏感度实测验证表}
\label{tab:sensitivity_verification}
\resizebox{\textwidth}{!}{%
\begin{tabular}{lcccc}
\toprule
\textbf{统计指标} & \textbf{代码声明敏感度 $\Delta_2$} & \textbf{最大经验位移 $\|f(D)-f(D')\|_2$} & \textbf{数学边界约束} & \textbf{审计结论} \\
\midrule
\textbf{User Prevalence} & $3.0000$ & $3.0000$ & $\le \sqrt{9} = 3.0$ & \textbf{PASS} (严格等于紧致界) \\
\textbf{Normalized Share} & $1.0000$ & $1.0000$ & $\le 1.0$ & \textbf{PASS} (严格受限于单纯形) \\
\textbf{Interaction Volume} & $20.0000$ & $20.0000$ & $\le C = 20.0$ & \textbf{PASS} (截断边界有效) \\
\textbf{Success Rate} & $28.2843$ & $20.0728$ & $\le \sqrt{2}C_{\text{rate}} = 28.2843$ & \textbf{PASS} (完全处于理论界内) \\
\textbf{Top-K Tasks} & $2.2361$ & $1.7321$ & $\le \sqrt{3} \approx 1.7321$ & \textbf{PASS} (受限于字典大小) \\
\bottomrule
\end{tabular}%
}
\end{table}

实测证明，所有机制在各类极端输入下均无一例外受控于声明敏感度，未发现任何敏感度低估或边界溢出漏洞。

\subsection{高斯采样器蒙特卡洛统计检验 (Monte Carlo Check)}
为了验证底层高斯采样器与解析参数校准的统计一致性，系统在基准配置（$\sigma = 12.6740$）下开展了 5,000 次独立高斯采样蒙特卡洛试验：
\begin{itemize}[leftmargin=2em]
  \item \textbf{均值收敛性检验：} 各单元经验均值在 $[-0.3339, +0.4354]$ 之间波动，最大均值偏差 $Z\text{-score} = 2.43 < 3.0$（符合零均值正态假设）；
  \item \textbf{标准差吻合性检验：} 各单元经验标准差在 $[12.3997, 12.9112]$ 之间，相对理论标度 $\sigma = 12.6740$ 的最大相对误差仅为 $\mathbf{2.16\% < 3.0\%}$。
\end{itemize}
这排除了实现中因随机数尺度换算错误导致的隐私保证失效。

\subsection{RDP 记账公式高精度理论对齐}
针对 12 个默认评估阶数 $\alpha \in [1.25, 128.0]$，测试程序对比了理论公式 $\rho(\alpha) = \frac{\alpha \Delta_2^2}{2\sigma^2}$ 与底层高精度 \texttt{Decimal} 记账模块的输出值：
\begin{itemize}[leftmargin=2em]
  \item 单次高斯释放的理论散度与代码计算值之间的绝对差值最大仅为 $\mathbf{4.44 \times 10^{-16}}$；
  \item 连续 5 次机制组合后的累积误差最大仅为 $\mathbf{2.22 \times 10^{-16}}$。
\end{itemize}
验证了记账引擎不仅在数学公式上完全正确，而且借助高精度定点数实现，有效规避了浮点数累加导致的有效数字丢失风险。

\subsection{规范重放、作用域隔离与二叉树节点复用验证}
在隔离状态测试中，验证了以下关键工程边界：
\begin{enumerate}[label=(\alph*), leftmargin=2.5em]
  \item \textbf{规范重放零消耗：} 对同一查询执行 Replay，增量 $\Delta\varepsilon = 0.0000$，\texttt{cache\_hit = True}，返回结果与首次发布完全一致；
  \item \textbf{外部多租户预算连通：} 供应商 \texttt{provider\_a} 发布后，\texttt{provider\_b} 在 \texttt{GLOBAL\_}\allowbreak\texttt{EXTERNAL} 共享作用域下立即同步累积消耗（$1.3421$），合谋防护生效；
  \item \textbf{作用域内二叉树节点按位恒等复用：} 在同一作用域（如 \texttt{GLOBAL\_EXTERNAL}）下，多个不同接收方查询重叠时间段时，系统从持久化数据库提取的相同规范节点加噪向量在浮点表示上\textbf{按位完全相等（Bitwise Identical）}，正常服务重启后依然保持确定性一致；而在独立作用域之间严格遵循隔离与记账规范。
\end{enumerate}

\section{系统局限性与后续技术演进路线}
\label{sec:limitations}

\subsection{交互总量高流量截断偏差及其优化方向}
在当前基准测试中，交互总量（Interaction Volume）的全局截断阈值固定设置为 $C=20$。由于真实企业环境存在大量依赖自动化工具或深度集成 IDE 的高频研发人员，固定截断导致整体交互总量呈现约 $39.8\%$ 的结构性低估（Under-counting Bias）。

后续技术攻关的核心在于研究\textbf{截断阈值 $C$ 与隐私参数 $\varepsilon$ 的联合优化模型（$C \times \varepsilon$ Joint Calibration）}：
\begin{equation}
  \min_{C} \mathbb{E}\left[ \left( \hat{M}(C) - M \right)^2 \right] = \min_{C} \left[ \text{Bias}^2(C) + \text{Var}(C, \varepsilon) \right]
\end{equation}
即在“截断造成的负向偏差”与“增大 $C$ 导致敏感度上升从而放大的高斯方差”之间寻找理论最优平衡点。

\subsection{自适应贡献截断的隐私理论考量}
针对截断阈值 $C$ 的选择，一种工程直觉是采用“基于数据分布自适应选择分位数（如 90\% 分位数）”。然而，在严格的差分隐私理论下必须清醒认识到：\textbf{直接从待保护的数据集自身计算分位数并将阈值反馈用于后续截断，该过程本身构成了对私有数据的观测，若不分配隐私预算将产生严重的隐私泄露}。

因此，后续关于自适应截断（Adaptive Clipping）的演进路线规划如下：
\begin{enumerate}[label=(\alph*), leftmargin=2.5em]
  \item \textbf{独立先验校准集方案：} 利用历史已过期、不属于当前敏感窗口的公共或历史交互分布估算相对稳定的经验阈值 $C$；
  \item \textbf{差分隐私私有分位数分配方案：} 将单次查询的总预算拆分为两部分：分配 $\varepsilon_{\text{clip}}$ 运行满足差分隐私的私有分位数估计机制，求出加噪截断阈值 $\tilde{C}$；随后分配 $\varepsilon_{\text{sum}} = \varepsilon - \varepsilon_{\text{clip}}$ 执行截断后的高斯聚合。该方案仍需在效用收益与分段预算损耗之间开展系统性实证评估。
\end{enumerate}

\subsection{细粒度业务分类体系的拓展适配}
当前系统的基准分类体系包含 3 个核心任务（Technical、Document、Analysis）。在企业实际生产场景中，业务部门往往期望进一步拆分为数十个子类别（例如代码审查、单元测试生成、SQL 编写、长文润色、合规审查等）。

随着分类维度 $D$ 的扩张，稀疏单元（Sparsity）数量显著增加。在微小调用单元中，高斯机制的随机扰动相对占比将增大。后续研发将探索树状层级分类法（Hierarchical Taxonomy）与分层差分隐私投影算法，在宏观层级保持高精度的同时，在微观叶子层级实施自适应样本抑制。

\subsection{更多业务维度的扩展规划}
在当前五类核心统计指标的基础上，后续版本将逐步演进支持更多维度的指标 Profile：
\begin{itemize}[leftmargin=2em]
  \item \textbf{Token 吞吐量统计（Token Volume Profile）：} 针对单次交互中数千乃至数万的输入/输出 Token 量，推导基于重尾对数截断的差分隐私释放方案；
  \item \textbf{推理延迟与 SLA 达标率（Latency Profile）：} 评估不同供应商在峰值负载下的服务响应时间分布；
  \item \textbf{错误与降级重试率（Fallback Rate Profile）：} 评估模型网关主备切换与重试频次，指导高可用架构调优。
\end{itemize}

\subsection{智能化效用规划器 (Utility-Aware Planner) 的演进}
当前版本已实现了基础的离线效用对比、误差拆解与 $\varepsilon$-Sweep 模拟。下一阶段的研发目标是将规划器（Planner）从“离线辅助分析”演进为“智能化闭环推荐”：即系统根据输入数据集规模 $N$、历史活跃特征以及目标接收方的剩余隐私配额，自动为用户推荐兼顾合规底线与数据可用性的候选发布时间窗口、推荐 $\varepsilon$ 建议值以及配套的截断参数建议。


% 参考文献
\newpage
\bibliographystyle{plain}
\bibliography{references}

% 附录
\newpage
\appendix
\section{核心数学定理推导与形式化证明}
\label{app:math}

\subsection{用户使用率 (Prevalence) 紧致敏感度证明}
\begin{techbox}{定理 A.1 (Prevalence $L_2$ 敏感度界限)}
设企业分类单元总维度为 $D$，单用户参与单元数受限于 $\|b_u\|_0 \le k$。则对直方图聚合查询 $H(D) = \sum_{u \in \mathcal{U}} b_u$，在用户级加/删邻接关系下的全局 $L_2$ 敏感度为 $\Delta_2 = \sqrt{\min(k, D)}$。
\end{techbox}

\begin{proof}
设相邻数据集为 $D_1$ 与 $D_2$，且 $D_2 = D_1 \cup \{u^*\}$（增加一个用户 $u^*$）。则聚合直方图的差值为：
\begin{equation}
  H(D_2) - H(D_1) = \sum_{u \in D_2} b_u - \sum_{u \in D_1} b_u = b_{u^*}
\end{equation}
根据欧几里得范数的定义：
\begin{equation}
  \|H(D_2) - H(D_1)\|_2 = \|b_{u^*}\|_2 = \sqrt{\sum_{j=1}^D b_{u^*, j}^2}
\end{equation}
由于每个单元的指示变量 $b_{u^*, j} \in \{0, 1\}$，故 $b_{u^*, j}^2 = b_{u^*, j}$。因此：
\begin{equation}
  \sum_{j=1}^D b_{u^*, j}^2 = \sum_{j=1}^D b_{u^*, j} = \|b_{u^*}\|_0
\end{equation}
由截断算子的定义，非零单元数不得超过 $k$；同时，非零单元数绝不可能超过系统总维度 $D$。因此 $\|b_{u^*}\|_0 \le \min(k, D)$。代入得：
\begin{equation}
  \|H(D_2) - H(D_1)\|_2 \le \sqrt{\min(k, D)}
\end{equation}
当构造一个在 $\min(k, D)$ 个单元中每个单元均恰好调用 1 次的用户时，等号严格成立。证毕。
\end{proof}

\subsection{用户归一化占比 (Normalized Share) 天然有界性证明}
\begin{techbox}{定理 A.2 (Normalized Share 敏感度上界)}
设用户工作流占比向量定义为 $x_u = m_u / \|m_u\|_1$（若 $\|m_u\|_1 = 0$ 则为 $\mathbf{0}$）。则聚合查询 $X(D) = \sum_{u \in \mathcal{U}} x_u$ 的全局 $L_2$ 敏感度满足 $\Delta_2 \le 1.0$。
\end{techbox}

\begin{proof}
对于任意被添加或删除的员工个体 $u^*$：
若 $\|m_{u^*}\|_1 = 0$，则 $x_{u^*} = \mathbf{0}$，$\|x_{u^*}\|_2 = 0 \le 1.0$。
若 $\|m_{u^*}\|_1 > 0$，向量 $x_{u^*}$ 各分量满足 $x_{u^*, j} \ge 0$ 且 $\sum_{j=1}^D x_{u^*, j} = 1$。展开 $L_2$ 范数的平方：
\begin{equation}
  \|x_{u^*}\|_2^2 = \sum_{j=1}^D x_{u^*, j}^2 \le \sum_{j=1}^D x_{u^*, j}^2 + \sum_{i \ne j} x_{u^*, i} x_{u^*, j} = \left(\sum_{j=1}^D x_{u^*, j}\right)^2 = \|x_{u^*}\|_1^2 = 1.0
\end{equation}
两侧开平方可得 $\|x_{u^*}\|_2 \le 1.0$。因此：
\begin{equation}
  \Delta_2 = \max_{u^*} \|X(D \cup \{u^*\}) - X(D)\|_2 = \max_{u^*} \|x_{u^*}\|_2 \le 1.0
\end{equation}
当某个用户的所有调用全部集中于单一单元时，$\|x_{u^*}\|_2 = 1.0$ 达到上界。证毕。
\end{proof}

\subsection{任务成功率联合向量敏感度证明}
\begin{techbox}{定理 A.3 (Success Rate 联合向量敏感度界限)}
设用户成功量向量为 $s_u$，总量向量为 $n_u$，满足 $0 \le s_{u, j} \le n_{u, j}$。在共享截断算子 $a_u = \min(1, C_{\text{rate}}/\|n_u\|_2)$ 下，联合向量 $w_u = [a_u s_u ; a_u n_u] \in \mathbb{R}^{2D}$ 的全局 $L_2$ 敏感度满足 $\Delta_2 \le \sqrt{2} C_{\text{rate}}$。
\end{techbox}

\begin{proof}
对于任意新增或删除的员工 $u^*$：
\begin{equation}
  \|w_{u^*}\|_2^2 = \|\bar{s}_{u^*}\|_2^2 + \|\bar{n}_{u^*}\|_2^2 = a_{u^*}^2 \|s_{u^*}\|_2^2 + a_{u^*}^2 \|n_{u^*}\|_2^2 = a_{u^*}^2 \left( \sum_{j=1}^D s_{u^*, j}^2 + \sum_{j=1}^D n_{u^*, j}^2 \right)
\end{equation}
由于单次调用的成功数不可能超过调用总数，故对于每一个分量均有 $0 \le s_{u^*, j} \le n_{u^*, j}$，从而 $s_{u^*, j}^2 \le n_{u^*, j}^2$。因此：
\begin{equation}
  \sum_{j=1}^D s_{u^*, j}^2 \le \sum_{j=1}^D n_{u^*, j}^2 = \|n_{u^*}\|_2^2
\end{equation}
将其代入联合范数不等式中：
\begin{equation}
  \|w_{u^*}\|_2^2 \le a_{u^*}^2 \left( \|n_{u^*}\|_2^2 + \|n_{u^*}\|_2^2 \right) = 2 a_{u^*}^2 \|n_{u^*}\|_2^2
\end{equation}
由截断比例因子定义 $a_{u^*} = \min\left(1, \frac{C_{\text{rate}}}{\|n_{u^*}\|_2}\right)$，若 $\|n_{u^*}\|_2 \le C_{\text{rate}}$，则 $a_{u^*} = 1$，此时 $a_{u^*} \|n_{u^*}\|_2 \le C_{\text{rate}}$；若 $\|n_{u^*}\|_2 > C_{\text{rate}}$，则 $a_{u^*} \|n_{u^*}\|_2 = C_{\text{rate}}$。在所有情形下均有：
\begin{equation}
  a_{u^*} \|n_{u^*}\|_2 \le C_{\text{rate}}
\end{equation}
代入可得 $\|w_{u^*}\|_2^2 \le 2 C_{\text{rate}}^2$，两侧开方即得：
\begin{equation}
  \|w_{u^*}\|_2 \le \sqrt{2} C_{\text{rate}}
\end{equation}
因此联合查询的全局敏感度满足 $\Delta_2 \le \sqrt{2} C_{\text{rate}}$。证毕。
\end{proof}

\section{工程验证与复现执行指南}
\label{app:tests}

本附录提供 \textbf{Private LLM Analytics} 系统的完整工程验证、测试执行与实验复现指令集。所有测试与实验均基于统一的沙盒与环境隔离规范设计，旨在为独立评审人提供确定性、高置信度的端到端复验路径。

\subsection{环境准备与依赖安装}

系统要求 Python 3.10 及以上运行环境（已验证 Python 3.10、3.11 及 3.12 兼容性）。建议在全新的虚拟环境中完成初始化配置：

\begin{techbox}{Python 虚拟环境初始化指令}
\begin{verbatim}
# 1. 创建独立虚拟环境
python -m venv .venv

# 2. 激活虚拟环境 (Windows PowerShell)
.\.venv\Scripts\Activate.ps1
# 或 Linux / macOS: source .venv/bin/activate

# 3. 升级核心包管理工具并安装项目依赖
python -m pip install --upgrade pip
pip install -r requirements.txt

# 4. 可选：以开发可编辑模式安装本核心库
pip install -e .
\end{verbatim}
\end{techbox}

核心运行依赖库包括：\texttt{numpy}、\texttt{scipy}、\texttt{pandas}、\texttt{pydantic}、\texttt{pyyaml}、\texttt{fastapi}、\texttt{uvicorn}、\texttt{jinja2}；测试框架采用 \texttt{pytest}。

\subsection{全量工程测试套件运行}

系统工程实现包含单元测试、机制边界测试、多工作线程记账并发测试、规范重放缓存测试以及时间二叉树持续观测测试等。执行全量测试套件：

\begin{techbox}{全量自动化回归测试指令}
\begin{verbatim}
# 运行全部 542 项工程测试套件
pytest tests -q

# 期望输出基准（耗时约 500-600 秒）：
# 542 passed, 0 failed in 573.10s
\end{verbatim}
\end{techbox}

若需针对特定子系统进行快速单模块审计，可运行下列分项测试指令：
\begin{itemize}[itemsep=2pt,topsep=2pt]
  \item \textbf{差分隐私与敏感度测试：} \verb|pytest tests/test_dp_mechanisms.py -v|
  \item \textbf{RDP 记账与多租户测试：} \verb|pytest tests/test_accounting_scopes.py -v|
  \item \textbf{规范重放与零预算穿透测试：} \verb|pytest tests/test_canonical_replay.py -v|
  \item \textbf{持续观测二叉树前缀测试：} \verb|pytest tests/test_continual_binary_tree.py -v|
  \item \textbf{Web 控制台与 API 接口测试：} \verb|pytest tests/test_web_console.py -v|
\end{itemize}

\subsection{全数值链端到端诊断审计}

为验证从原始遥测日志读取、用户贡献截断、全局敏感度校准、高斯噪声添加、非负保序后处理、RDP 记账扣除，到 API 封装与 CSV 导出的完整数值一致性，系统提供了统一数值诊断线束：

\begin{techbox}{端到端数值链审计执行}
\begin{verbatim}
python experiments/audit_dp_mechanisms.py
\end{verbatim}
\end{techbox}

该线束针对 $N=1,000$ 典型用户样本，在固定随机种子下遍历全部 5 项统计指标，执行下列严密检查：
\begin{enumerate}[itemsep=1pt,topsep=2pt]
  \item \textbf{原始精确值计算（Raw Exact）：} 验证无截断下的真实业务统计基准。
  \item \textbf{贡献限制校验（Bounded Exact）：} 验证截断后用户特征向量严格落在界限内（如 Prevalence $\le 9$、Volume $L_2 \le 20.0$）。
  \item \textbf{敏感度与方差推导：} 打印精确的 $\Delta_2$ 与由 Balle-Wang 解析算法求解出的最小方差 $\sigma$。
  \item \textbf{高斯噪声释放与后处理：} 验证加噪后向量经非负约束、加和归一化或阈值抑制后，返回的业务结果在语义范围内（如成功率 $\in [0, 1]$ 或 \texttt{SUPPRESSED}）。
  \item \textbf{规范重放判定：} 执行相同参数重放，验证服务端精准命中 \verb|cache_hit = True|，预算消耗 $\Delta\varepsilon = 0.0$。
  \item \textbf{接收方隔离与 CSV 导出：} 验证向外部接收方（如 \texttt{provider\_a}）请求时，返回结果已被自动裁剪且不可见其他提供商数据，同时校验生成的下载 CSV 文件格式与数值完全一致。
\end{enumerate}
运行结果最终输出各指标全数值链状态为 \verb|STATUS: PASS|。

\subsection{数学边界与蒙特卡洛独立校验}

为排除对第三方统计库的先验假设，运行独立深度数值校验脚本：

\begin{techbox}{数学敏感度与蒙特卡洛校验}
\begin{verbatim}
python scratch/audit_deep_checks.py
\end{verbatim}
\end{techbox}

该测试执行四项独立的底层校验：
\begin{itemize}[itemsep=2pt,topsep=2pt]
  \item \textbf{经验敏感度网格搜索：} 构建 6 类处于边界条件下的极端用户向量（全零调用、单单元脉冲式调用、均匀满负荷调用、极端倾斜调用），遍历所有相邻数据集对 $D \sim D'$，验证 $\|f(D) - f(D')\|_2 \le \Delta_2$ 严格成立，且在最劣测试用例下等号取到。
  \item \textbf{蒙特卡洛噪声校准验证：} 针对已求得的 $\sigma$，进行 5,000 次独立高斯采样，验证样本均值偏差在 $3\sigma/\sqrt{M}$ 统计置信区间内，样本标准差相对理论误差小于 $3.0\%$。
  \item \textbf{RDP 记账解析精度校验：} 将代码库计算的 Rényi 散度 $\rho(\alpha)$ 与解析理论值 $\frac{\alpha \Delta_2^2}{2\sigma^2}$ 进行逐阶对比，在 12 个评价阶数上均达到浮点机器精度级别（残差 $< 10^{-15}$）。
  \item \textbf{动态分类维度泛化校验：} 分别在 $2\times 2$、$3\times 3$ 以及 $4\times 5$ 的非对称分类维度下实例化全部指标算子，确认无静态维度硬编码缺陷。
\end{itemize}

\subsection{规范重放与持续观测一致性复现}

验证多租户环境下的预算账本隔离与时间二叉树节点重用特性：

\begin{techbox}{重放机制与二叉树节点复验}
\begin{verbatim}
python scratch/audit_replay_and_continual.py
\end{verbatim}
\end{techbox}

该脚本验证两项系统关键行为：
\begin{enumerate}[itemsep=1pt,topsep=2pt]
  \item \textbf{接收方作用域穿透验证：} 外部接收方 \verb|provider_a| 触发首次发布消耗 $\varepsilon=1.3421$ 后，同属全局外部作用域的 \verb|provider_b| 请求同一统计指标时，系统直接返回已有密态发布结果且不追加扣除外部预算账本；而内部受信任接收方 \verb|internal| 拥有独立的内部审计账本，互不挤占。
  \item \textbf{二叉树持久化与作用域节点复用确定性：} 验证跨进程重启后，二叉树在同一隐私作用域内对已完成发布的历史区间节点实现按位完全一致的重放复用输出，且仅在时间窗口拓展时对新激活的二叉树父节点进行有限增量记账。
\end{enumerate}

\subsection{本地可视化 Web 控制台交互启动}

系统提供了面向非技术利益相关者的可解释性 Web 控制台：

\begin{techbox}{启动本地可视化服务}
\begin{verbatim}
python -m private_llm_analytics.web.app
# 服务将监听并在本地启动：http://127.0.0.1:8000
\end{verbatim}
\end{techbox}

评审人员可在浏览器中打开 \texttt{http://127.0.0.1:8000}，体验端到端交互流：
\begin{itemize}[itemsep=1pt,topsep=2pt]
  \item \textbf{数据接入：} 下载示例 CSV 模板，上传企业遥测日志或系统预置的合成数据集。
  \item \textbf{统计问题选择：} 切换“用户采纳率”、“工作流占比”、“总调用规模”与“成功率”。
  \item \textbf{通俗化机制释义：} 控制台默认以“单人最大影响限制”、“防护噪声”与“非负约束”等业务语言展示，并提供“展开数学细节”选项支持实时 KaTeX 公式查看。
  \item \textbf{误差可视化对比：} 观察原始真实值（Raw Exact）、有界精确值（Bounded Exact）与加噪发布值（DP Estimate）之间的数值偏差。
  \item \textbf{合规结果导出：} 点击一键下载脱敏分析结果 CSV 文件，验证接收方作用域过滤效果。
\end{itemize}

\subsection{生产账本隔离与状态安全守则}

\begin{techbox}[colback=white,colframe=accentamber]{注意事项：持久化账本操作安全规范}
系统生产运行的默认账本存储于 \texttt{results/privacy\_state.db}。在进行工程复现与代码审计时，请严格遵守下列操作规范：
\begin{enumerate}[itemsep=1pt,topsep=1pt]
  \item \textbf{禁止手动篡改或重置生产账本：} \texttt{privacy\_state.db} 记录了已发布事件的历史序列与真实已耗预算，在生产环境下任何对数据库的清空或回滚均会导致严重的隐私泄露与噪声平均攻击（Noise Averaging Attack）。
  \item \textbf{测试脚本强制使用独立数据库：} 所有自动化测试与实验脚本均通过 \texttt{IsolatedStateDB} 上下文管理器在系统临时目录生成临时 SQLite 实例，并在测试生命周期结束后自动清理，避免测试副作用污染生产账本。
  \item \textbf{版本控制状态监控：} 在完成所有本地验证操作后，请执行 \texttt{git status}，确保 \texttt{results/privacy\_state.db} 保持未修改状态。
\end{enumerate}
\end{techbox}

\section{生产环境配置规范与 YAML 范例}
\label{app:config}

\textbf{Private LLM Analytics} 系统遵循声明式配置与代码逻辑严格解耦的软件工程原则。所有隐私界限、机制参数、记账策略、外部接收方权限与领域分类均通过标准 YAML 文件进行配置驱动。本附录提供生产环境的核心配置架构解析与基准模板。

\subsection{指标机制定义配置 (metrics.yaml)}

\texttt{configs/metrics.yaml} 统一定义系统中所有已注册统计指标的维度结构、用户级贡献限制模式与噪声释放机制类型。

\begin{techbox}{指标机制生产配置清单}
\begin{verbatim}
metrics:
  # 1. 用户使用率指标
  user_prevalence:
    enabled: true
    type: prevalence
    dimensions:
      - provider
      - task
    contribution:
      mode: binary              # 二值指示向量
      max_cells_per_user: 9     # 单用户最多参与单元数截断界限 k
    mechanism:
      type: gaussian

  # 2. 用户归一化使用占比指标
  user_normalized_share:
    enabled: true
    type: normalized_share
    dimensions:
      - provider
      - task
    contribution:
      mode: l1_normalized       # L1 归一化用户工作流向量
      norm_bound: 1.0           # 天然敏感度上界
    mechanism:
      type: gaussian

  # 3. 交互调用总量指标
  interaction_volume:
    enabled: true
    type: volume
    dimensions:
      - provider
      - task
    contribution:
      mode: l2_clip             # 用户多维调用频次 L2 范数截断
      clip_norm: 20.0           # 截断界限 C_vol
    mechanism:
      type: gaussian

  # 4. 任务成功率指标
  success_rate:
    enabled: true
    type: success_rate
    dimensions:
      - provider
      - task
    contribution:
      mode: joint_l2_clip       # 共享调用总量总范数的联合截断
      clip_norm_rate: 20.0      # 成功率截断界限 C_rate
      min_noisy_denominator: 1.0 # 有噪分母抑制最小阈值 tau
    mechanism:
      type: gaussian

  # 5. 热门任务发现指标
  top_k_tasks:
    enabled: true
    type: top_k_tasks
    dimensions:
      - task
    contribution:
      mode: binary
      max_tasks_per_user: 5
      top_k: 5                  # 最终加噪截断输出的前 K 个任务
    mechanism:
      type: gaussian
\end{verbatim}
\end{techbox}

\subsection{全生命周期隐私与规划器配置 (privacy.yaml)}

\texttt{configs/privacy.yaml} 集中配置全系统的差分隐私基础假设、生命周期预算上限、RDP 记账阶数、外部接收方合并模式以及数据驱动离线规划器（Offline Planner）的先验分布。

\begin{techbox}{隐私管理与记账生产配置清单}
\begin{verbatim}
privacy:
  privacy_unit: user            # 隐私保护实体：单个自然人用户
  adjacency: add_remove         # 邻接关系：加/删单用户所有记录

  horizon:
    days: 90                    # 规划与记账生命周期时间跨度

  global_budget:
    epsilon: 4.0                # 全局生命周期累计隐私预算上限 Epsilon
    delta: 4.0e-6               # 全局累计松弛项 Delta

  default_query_budget:
    epsilon: 1.0                # 单次常规发布默认隐私预算
    delta: 1.0e-6

  accountant:
    mode: rdp                   # 采用 Rényi 差分隐私紧致组合记账
    epsilon_cap: 4.0
    delta_global: 1.0e-6
    # 跟踪的 Rényi 散度阶数列表 alpha
    orders:
      - 1.25
      - 1.5
      - 2.0
      - 3.0
      - 4.0
      - 5.0
      - 8.0
      - 10.0
      - 16.0
      - 32.0
      - 64.0
      - 128.0

  history:
    mode: lifetime              # 永久追踪全部发布历史事件

  recipient_composition:
    mode: global_external       # 所有外部模型提供商共用同一外部预算池

  default_mechanism:
    type: gaussian              # 核心机制：解析高斯机制 (Balle-Wang)

  release:
    # 开启规范发布语义缓存，相同查询零预算开销
    reuse_existing_release: true
    allow_budget_overrun: false    # 预算耗尽时执行 Fail-Closed 熔断

  planner:
    enabled: true
    mode: static_public         # 基于公开先验参数的离线规划
    prior:
      distribution: contribution_model
      mean: 60.0
      std: 20.0
      lower: 5.0
      upper: 150.0
      public_seed: 900001
      num_reference_users: 2000
    weights:
      objective_mode: EXPECTED_RELEASE_L1
      data: 1.0
      clipping: 1.0
      noise: 1.0
      aggregation: 1.0
      staleness: 0.5
      window: 0.2
    candidate_clipping: [5.0, 10.0, 20.0, 40.0, 80.0]
    candidate_windows: [7, 14, 30, 60, 90]
    candidate_frequencies: [1, 7, 14, 30]
\end{verbatim}
\end{techbox}

\subsection{多接收方权限与作用域配置 (recipients.yaml)}

\texttt{configs/recipients.yaml} 实现了基于最小权限原则的强隔离策略。系统将接收方划分为内部受信任审计方与外部竞争性模型服务商，并对外部实体施加列级过滤与跨供应商视图屏蔽。

\begin{techbox}{接收方角色与权限矩阵配置}
\begin{verbatim}
recipients:
  # 内部受信任治理团队
  internal:
    role: trusted
    display_name: "内部 AI 治理团队"
    providers: ["*"]            # 可查看所有模型供应商
    metrics: ["*"]              # 可访问全部指标体系
    exact_results: true         # 允许在受控沙盒内访问精确原始统计
    dp_results: true
    cross_provider_comparison: true

  # 外部模型供应商 A
  provider_a:
    role: external
    display_name: "OpenAI"
    providers:
      - provider_a              # 仅允许查看其自身模型对应的统计切片
    metrics:
      - user_prevalence
      - user_normalized_share
      - interaction_volume
      - success_rate
    exact_results: false        # 严格禁止访问明细或精确基准值
    dp_results: true            # 仅允许接收差分隐私脱敏释放结果
    cross_provider_comparison: false # 强隔离：屏蔽所有横向对比字段

  # 外部模型供应商 B
  provider_b:
    role: external
    display_name: "Kimi"
    providers:
      - provider_b
    metrics:
      - user_prevalence
      - user_normalized_share
      - interaction_volume
      - success_rate
    exact_results: false
    dp_results: true
    cross_provider_comparison: false

  # 外部模型供应商 C
  provider_c:
    role: external
    display_name: "DeepSeek"
    providers:
      - provider_c
    metrics:
      - user_prevalence
      - user_normalized_share
      - interaction_volume
      - success_rate
    exact_results: false
    dp_results: true
    cross_provider_comparison: false
\end{verbatim}
\end{techbox}

\subsection{分类体系与基准场景配置 (taxonomy.yaml \& scenario.yaml)}

系统通过独立分类配置实现对多模型及业务任务的动态扩展能力，避免在代码内部硬编码业务标签。

\begin{techbox}{分类与基准场景配置}
\begin{verbatim}
# configs/taxonomy.yaml
taxonomy:
  version: 1
  ordering: "canonical_lexicographic" # 强制采用规范字典序对齐多维向量
  providers:
    - provider_a
    - provider_b
    - provider_c
  tasks:
    - technical
    - document
    - analysis

# configs/scenario.yaml
dataset:
  source: synthetic
population:
  mode: fixed_public_population
  n_users: 50000              # 企业级合成基准总体规模：50,000 名员工
  n_days: 30                  # 观测周期跨度：30 天
usage_frequency:
  distribution: normal
  mean: 60                    # 单用户平均调用频次基准
  std: 20
  minimum: 5
  maximum: 150
random_seed: 42
\end{verbatim}
\end{techbox}


\end{document}
